%%%%%%%%%%%%%%%%%%%%%%%%%%%%%%%%%%%%%%%%%%%%%%%%%%%%%%%%%%%%%%%%%%%%%%%%%%%%%%%%
\chapter{Alternative Stochastic Approaches to Quantum Theory}
\label{ch:AlternativeStochastic}
%%%%%%%%%%%%%%%%%%%%%%%%%%%%%%%%%%%%%%%%%%%%%%%%%%%%%%%%%%%%%%%%%%%%%%%%%%%%%%%%
%%%%%%%%%%%%%%%%%%%%%%%%%%%%%%%%%%%%%%%%%%%%%%%%%%%%%%%%%%%%%%%%%%%%%%%%%%%%%%%%
\section{Introduction}
%%%%%%%%%%%%%%%%%%%%%%%%%%%%%%%%%%%%%%%%%%%%%%%%%%%%%%%%%%%%%%%%%%%%%%%%%%%%%%%%
The remarkable predictive success of quantum mechanics has never
eliminated interest in the possibility that it may emerge from a deeper,
more fundamental theory. Since the earliest days of quantum mechanics,
many physicists have questioned whether the probabilistic character of
the theory is truly fundamental or merely reflects our ignorance of
underlying microscopic dynamics.
One possible answer is provided by hidden-variable theories, of which
the de Broglie--Bohm theory is the best-known example. Another,
conceptually different, possibility is that quantum behavior arises from
an underlying stochastic process. In this picture the apparent
indeterminacy of quantum mechanics does not originate from intrinsic
probability, but from continual interaction with microscopic random
degrees of freedom.
The idea has appeared repeatedly during the history of physics in
several independent forms. Although these approaches differ
substantially in their mathematical formulation, they all attempt to
answer essentially the same question:
\begin{quote}
\emph{Can quantum mechanics be understood as the statistical description
of an underlying stochastic dynamics?}
\end{quote}
This question has motivated a surprisingly diverse collection of
theories.
Some authors assume that particles themselves execute Brownian motion.
Others attribute the stochasticity to classical electromagnetic fields,
to random space-time geometry, to non-differentiable trajectories, to
matrix-valued microscopic dynamics, or even to stochastic evolution in
an auxiliary mathematical time.
Consequently, the term \emph{stochastic quantum theory} does not denote
a single well-defined framework but rather an entire family of
approaches sharing a common philosophical motivation while differing
greatly in physical assumptions.
The purpose of the present chapter is not merely to review these
theories individually. Equally important is to compare them
systematically, identify their common conceptual structure, and clarify
the physical questions they attempt to answer.
%%%%%%%%%%%%%%%%%%%%%%%%%%%%%%%%%%%%%%%%%%%%%%%%%%%%%%%%%%%%%%%%%%%%%%%%%%%%%%%%
\subsection{What Can Be Stochastic?}
%%%%%%%%%%%%%%%%%%%%%%%%%%%%%%%%%%%%%%%%%%%%%%%%%%%%%%%%%%%%%%%%%%%%%%%%%%%%%%%%
One of the principal sources of confusion in the literature is that the
word \emph{stochastic} is used in several fundamentally different ways.
The first possibility is that the particle trajectory itself is random.
This viewpoint is adopted in Nelson's stochastic mechanics.
A second possibility is that the particle remains completely classical,
while the electromagnetic field possesses a universal stochastic
background. This is the central hypothesis of stochastic
electrodynamics.
A third possibility is that the stochastic process takes place not in
physical space-time but in the infinite-dimensional space of field
configurations. This idea forms the basis of Parisi--Wu stochastic
quantization.
Other approaches attribute randomness to the geometry of space-time
itself, to fractal properties of trajectories, or to the statistical
mechanics of noncommuting matrices.
Thus, before comparing the various theories, one must first answer a
more fundamental question:
\begin{center}
\emph{What physical quantity is assumed to fluctuate?}
\end{center}
As we shall see, the answer to this question largely determines the
mathematical structure and physical interpretation of each theory.
%%%%%%%%%%%%%%%%%%%%%%%%%%%%%%%%%%%%%%%%%%%%%%%%%%%%%%%%%%%%%%%%%%%%%%%%%%%%%%%%
\subsection{Common Objectives}
%%%%%%%%%%%%%%%%%%%%%%%%%%%%%%%%%%%%%%%%%%%%%%%%%%%%%%%%%%%%%%%%%%%%%%%%%%%%%%%%
Despite their diversity, nearly all stochastic approaches pursue several
closely related goals.
These include
\begin{itemize}
\item explaining the origin of Planck's constant,
\item deriving the Schrödinger equation from more primitive dynamical
principles,
\item understanding the physical meaning of the wave function,
\item explaining quantum probabilities,
\item clarifying the measurement problem,
\item establishing a deeper connection between classical and quantum
physics.
\end{itemize}
Not every theory attempts to solve all of these problems.
For example, stochastic quantization is primarily a computational method
for constructing quantum field theories rather than an interpretation of
quantum mechanics.
Conversely, stochastic electrodynamics explicitly seeks a classical
physical origin of quantum phenomena.
Recognizing these different objectives is essential for a fair
comparison.
%%%%%%%%%%%%%%%%%%%%%%%%%%%%%%%%%%%%%%%%%%%%%%%%%%%%%%%%%%%%%%%%%%%%%%%%%%%%%%%%
\subsection{Criteria for Comparison}
%%%%%%%%%%%%%%%%%%%%%%%%%%%%%%%%%%%%%%%%%%%%%%%%%%%%%%%%%%%%%%%%%%%%%%%%%%%%%%%%
Throughout this chapter we shall compare the various theories according
to a common set of questions.
For each approach we shall ask:
\begin{enumerate}
\item What is the fundamental stochastic variable?
\item What is the evolution parameter?
\item Is the theory relativistic?
\item Does the theory explain the origin of Planck's constant?
\item Can it derive the Schrödinger equation?
\item Can it describe spin?
\item Does it naturally generalize to quantum field theory?
\item How are gauge interactions incorporated?
\item Is gravity included?
\item What are its principal successes?
\item What are its most serious unresolved problems?
\end{enumerate}
Using identical criteria for every theory allows a meaningful comparison
that is largely absent from the existing literature.
%%%%%%%%%%%%%%%%%%%%%%%%%%%%%%%%%%%%%%%%%%%%%%%%%%%%%%%%%%%%%%%%%%%%%%%%%%%%%%%%
\subsection{Organization of This Chapter}
%%%%%%%%%%%%%%%%%%%%%%%%%%%%%%%%%%%%%%%%%%%%%%%%%%%%%%%%%%%%%%%%%%%%%%%%%%%%%%%%
The chapter begins with stochastic electrodynamics, the oldest and most
extensively developed classical stochastic approach to quantum theory.
We then review stochastic quantization, scale relativity, trace
dynamics, and several more recent proposals.
The chapter concludes with a systematic comparison of all approaches and
a discussion of the lessons they provide for the stochastic
Hamiltonian framework developed in the remainder of this monograph.
%%%%%%%%%%%%%%%%%%%%%%%%%%%%%%%%%%%%%%%%%%%%%%%%%%%%%%%%%%%%%%%%%%%%%%%%%%%%%%%%
\subsection*{Editorial Comment}
%%%%%%%%%%%%%%%%%%%%%%%%%%%%%%%%%%%%%%%%%%%%%%%%%%%%%%%%%%%%%%%%%%%%%%%%%%%%%%%%
The diversity of stochastic approaches reviewed in this chapter should
not be regarded as evidence of failure. Rather, it reflects the
difficulty of the underlying problem.
Quantum mechanics is extraordinarily successful experimentally, and any
attempt to derive it from deeper principles faces exceptionally
stringent constraints. The existence of multiple stochastic programs,
developed largely independently over more than half a century,
demonstrates that the possibility of an underlying stochastic dynamics
has remained scientifically compelling despite the absence of a
universally accepted theory.
Our objective is therefore not to advocate a particular approach, but to
understand what each contributes, where each succeeds, where each fails,
and which ideas appear sufficiently robust to serve as ingredients of a
more complete relativistic stochastic theory.
%%%%%%%%%%%%%%%%%%%%%%%%%%%%%%%%%%%%%%%%%%%%%%%%%%%%%%%%%%%%%%%%%%%%%%%%%%%%%%%%
\section{Historical Development of Stochastic Electrodynamics}
\label{sec:SEDHistory}
%%%%%%%%%%%%%%%%%%%%%%%%%%%%%%%%%%%%%%%%%%%%%%%%%%%%%%%%%%%%%%%%%%%%%%%%%%%%%%%%
Stochastic electrodynamics (SED) is one of the oldest attempts to
understand quantum phenomena in terms of an underlying classical
stochastic dynamics. Unlike stochastic mechanics, where randomness is
attributed to the motion of particles, SED assumes that the
electromagnetic field itself possesses an intrinsic random component
which exists even at absolute zero temperature.
From this viewpoint particles always obey classical equations of motion.
Their apparently quantum behavior results from continual interaction
with an omnipresent stochastic electromagnetic background.
The central hypothesis may therefore be summarized in a single sentence:
\begin{quote}
Quantum phenomena originate not from quantization of matter, but from
the interaction of classical matter with a universal classical
stochastic radiation field.
\end{quote}
This idea has a long and fascinating history extending over more than a
century.
%%%%%%%%%%%%%%%%%%%%%%%%%%%%%%%%%%%%%%%%%%%%%%%%%%%%%%%%%%%%%%%%%%%%%%%%%%%%%%%%
\subsection{Origins in Classical Radiation Theory}
%%%%%%%%%%%%%%%%%%%%%%%%%%%%%%%%%%%%%%%%%%%%%%%%%%%%%%%%%%%%%%%%%%%%%%%%%%%%%%%%
The roots of SED lie in the crisis of classical radiation theory at the
end of the nineteenth century.
Maxwell's equations had achieved spectacular success in describing
electromagnetic phenomena. At the same time, classical statistical
mechanics appeared capable of explaining the thermodynamic properties of
matter.
The difficulty arose when these two theories were combined to describe
thermal radiation.
According to classical statistical mechanics, every normal mode of the
electromagnetic field should possess the average energy
\[\langle E\rangle=k_BT.\]
Since the number of modes increases rapidly with frequency, the total
energy diverges,
\[U=\int_0^\infty u(\omega)d\omega=\infty,\]
producing the well-known ultraviolet catastrophe.
Planck resolved this difficulty in 1900 by introducing energy quanta,
\[E=n\hbar\omega.\]
Although enormously successful, this proposal raised a profound
question.
Is quantization truly fundamental, or does it merely provide an
effective description of some deeper classical dynamics?
This question eventually became one of the principal motivations for
stochastic electrodynamics.
%%%%%%%%%%%%%%%%%%%%%%%%%%%%%%%%%%%%%%%%%%%%%%%%%%%%%%%%%%%%%%%%%%%%%%%%%%%%%%%%
\subsection{Einstein, Lorentz and Classical Fluctuations}
%%%%%%%%%%%%%%%%%%%%%%%%%%%%%%%%%%%%%%%%%%%%%%%%%%%%%%%%%%%%%%%%%%%%%%%%%%%%%%%%
Long before the modern formulation of SED, several pioneers had already
recognized the importance of fluctuations in electromagnetic theory.
Einstein's work on Brownian motion demonstrated that random microscopic
forces could produce observable macroscopic phenomena.
Lorentz investigated fluctuations generated by classical
electromagnetic fields.
Planck himself repeatedly reconsidered the possibility that the vacuum
might possess an intrinsic electromagnetic energy even at zero
temperature.
Although none of these ideas constituted stochastic electrodynamics in
its modern form, they established the conceptual foundations upon which
the later theory was built.
%%%%%%%%%%%%%%%%%%%%%%%%%%%%%%%%%%%%%%%%%%%%%%%%%%%%%%%%%%%%%%%%%%%%%%%%%%%%%%%%
\subsection{The Modern Beginning}
%%%%%%%%%%%%%%%%%%%%%%%%%%%%%%%%%%%%%%%%%%%%%%%%%%%%%%%%%%%%%%%%%%%%%%%%%%%%%%%%
Modern stochastic electrodynamics emerged during the 1960s.
The decisive step was the realization that Maxwell's equations admit a
Lorentz-invariant random solution whose spectral density is proportional
to
\[\omega^3.\]
Instead of quantizing the electromagnetic field, one assumes that this
random radiation field exists as a classical solution of Maxwell's
equations.
Among the principal contributors were
\begin{itemize}
\item Timothy H. Boyer,
\item Trevor W. Marshall,
\item Luis de la Peña,
\item Ana María Cetto,
\item Francisco Francia,
\item Theodore Puthoff,
\item Daniel Cole,
\item Yorick Savoy,
\item Gerhard Hegerfeldt (critical analyses),
\item and many others.
\end{itemize}
These authors transformed SED from an isolated idea into a coherent
research program.
%%%%%%%%%%%%%%%%%%%%%%%%%%%%%%%%%%%%%%%%%%%%%%%%%%%%%%%%%%%%%%%%%%%%%%%%%%%%%%%%
\subsection{The Boyer Revolution}
%%%%%%%%%%%%%%%%%%%%%%%%%%%%%%%%%%%%%%%%%%%%%%%%%%%%%%%%%%%%%%%%%%%%%%%%%%%%%%%%
The work of Boyer occupies a special position in the history of SED.
He demonstrated that many phenomena usually regarded as intrinsically
quantum could be reproduced by combining classical electrodynamics with
a Lorentz-invariant stochastic zero-point radiation field.
Among his best-known achievements are
\begin{itemize}
\item derivation of the Planck spectrum,
\item blackbody radiation,
\item harmonic oscillator,
\item Casimir forces,
\item diamagnetism,
\item thermal equilibrium,
\item Unruh-type phenomena.
\end{itemize}
Whether these results constitute a genuine derivation of quantum
mechanics remains controversial.
Nevertheless, Boyer's work demonstrated that classical physics is far
more powerful than previously believed once stochastic vacuum radiation
is included.
%%%%%%%%%%%%%%%%%%%%%%%%%%%%%%%%%%%%%%%%%%%%%%%%%%%%%%%%%%%%%%%%%%%%%%%%%%%%%%%%
\subsection{The Mexican School}
%%%%%%%%%%%%%%%%%%%%%%%%%%%%%%%%%%%%%%%%%%%%%%%%%%%%%%%%%%%%%%%%%%%%%%%%%%%%%%%%
A second major center of SED developed around Luis de la Peña and Ana
María Cetto.
Their approach emphasized the statistical mechanics of particles
interacting continuously with the zero-point field.
Particular attention was devoted to
\begin{itemize}
\item derivation of the Schrödinger equation,
\item interpretation of the wave function,
\item momentum fluctuations,
\item diffusion processes,
\item relation to Nelson's mechanics,
\item foundations of quantum mechanics.
\end{itemize}
Their comprehensive monograph remains one of the principal references in
the field.
%%%%%%%%%%%%%%%%%%%%%%%%%%%%%%%%%%%%%%%%%%%%%%%%%%%%%%%%%%%%%%%%%%%%%%%%%%%%%%%%
\subsection{Revival During the Twenty-First Century}
%%%%%%%%%%%%%%%%%%%%%%%%%%%%%%%%%%%%%%%%%%%%%%%%%%%%%%%%%%%%%%%%%%%%%%%%%%%%%%%%
Interest in SED increased again after approximately 2000 for several
reasons.
First,
high-precision numerical simulations became possible.
Second,
renewed interest in quantum foundations encouraged reconsideration of
alternative formulations of quantum mechanics.
Third,
connections emerged between SED and several modern topics including
\begin{itemize}
\item Casimir physics,
\item vacuum fluctuations,
\item nonequilibrium statistical mechanics,
\item quantum thermodynamics,
\item relativistic stochastic dynamics.
\end{itemize}
Although SED remains outside the mainstream interpretation of quantum
mechanics, it continues to generate active research.
%%%%%%%%%%%%%%%%%%%%%%%%%%%%%%%%%%%%%%%%%%%%%%%%%%%%%%%%%%%%%%%%%%%%%%%%%%%%%%%%
\subsection{Relation to Other Stochastic Theories}
%%%%%%%%%%%%%%%%%%%%%%%%%%%%%%%%%%%%%%%%%%%%%%%%%%%%%%%%%%%%%%%%%%%%%%%%%%%%%%%%
Historically, SED developed almost independently of Nelson's stochastic
mechanics.
This independence is rather surprising because both theories seek a
classical stochastic origin of quantum behavior.
The principal distinction lies in the source of stochasticity.
Nelson assumes
\[\boxed{\text{Particle motion}\longrightarrow\text{stochastic}.}\]
SED assumes
\[\boxed{\text{Electromagnetic field}\longrightarrow\text{stochastic}.}\]
This apparently simple difference has profound consequences for the
mathematical structure of the resulting theories.
%%%%%%%%%%%%%%%%%%%%%%%%%%%%%%%%%%%%%%%%%%%%%%%%%%%%%%%%%%%%%%%%%%%%%%%%%%%%%%%%
\subsection*{Critical Assessment}
%%%%%%%%%%%%%%%%%%%%%%%%%%%%%%%%%%%%%%%%%%%%%%%%%%%%%%%%%%%%%%%%%%%%%%%%%%%%%%%%
The historical development of stochastic electrodynamics illustrates an
important lesson.
SED should not be viewed merely as an unsuccessful attempt to replace
quantum mechanics.
Rather, it represents a sustained effort to determine how much of
quantum physics can be recovered from classical electrodynamics when the
vacuum is regarded as a genuine stochastic physical medium.
Whether or not this program ultimately succeeds is less important than
the insight it provides into the relationship between classical field
theory, stochastic processes and quantum phenomena.
From the viewpoint of the present monograph, SED constitutes one of the
principal alternative stochastic programs that any new relativistic
stochastic theory must carefully compare itself with.
%%%%%%%%%%%%%%%%%%%%%%%%%%%%%%%%%%%%%%%%%%%%%%%%%%%%%%%%%%%%%%%%%%%%%%%%%%%%%%%%
\section{Mathematical Foundations of Classical Zero-Point Radiation}
\label{sec:SEDMathematical}
%%%%%%%%%%%%%%%%%%%%%%%%%%%%%%%%%%%%%%%%%%%%%%%%%%%%%%%%%%%%%%%%%%%%%%%%%%%%%%%%
The defining postulate of stochastic electrodynamics is the existence of
a classical random electromagnetic field that persists even at absolute
zero temperature. Unlike ordinary thermal radiation, this field is
assumed to be a permanent property of the physical vacuum.
The particles themselves remain entirely classical. Their stochastic
motion arises solely through interaction with this universal
electromagnetic background.
Consequently, the mathematical formulation of SED begins not with
particle dynamics but with the statistical properties of solutions of
Maxwell's equations.
%%%%%%%%%%%%%%%%%%%%%%%%%%%%%%%%%%%%%%%%%%%%%%%%%%%%%%%%%%%%%%%%%%%%%%%%%%%%%%%%
\subsection{Free Electromagnetic Field}
%%%%%%%%%%%%%%%%%%%%%%%%%%%%%%%%%%%%%%%%%%%%%%%%%%%%%%%%%%%%%%%%%%%%%%%%%%%%%%%%
In the absence of charges and currents the Maxwell equations are
\begin{align}
\nabla\cdot\mathbf E &=0,
    &\nabla\cdot\mathbf B &=0,\\
\nabla\times\mathbf E&=-\frac1c\frac{\partial\mathbf B}{\partial t},
&\nabla\times\mathbf B&=\frac1c\frac{\partial\mathbf E}{\partial t}.
\end{align}
Introducing the four-potential
\[A^\mu=(\phi,\mathbf A),\]
and adopting the Lorenz gauge,
\[\partial_\mu A^\mu=0,\]
the field equations reduce to
\begin{equation}
\square A^\mu=0,
\label{eq:WaveEquation}
\end{equation}
where
\[\square=\eta^{\mu\nu}\partial_\mu\partial_\nu
=\frac1{c^2}\frac{\partial^2}{\partial t^2}-\nabla^2.\]
Equation (\ref{eq:WaveEquation}) possesses infinitely many plane-wave
solutions.
SED assumes that the physical vacuum contains a random superposition of
all such modes.
%%%%%%%%%%%%%%%%%%%%%%%%%%%%%%%%%%%%%%%%%%%%%%%%%%%%%%%%%%%%%%%%%%%%%%%%%%%%%%%%
\subsection{Plane-Wave Expansion}
%%%%%%%%%%%%%%%%%%%%%%%%%%%%%%%%%%%%%%%%%%%%%%%%%%%%%%%%%%%%%%%%%%%%%%%%%%%%%%%%
The vector potential is expanded as
\begin{equation}
\mathbf A(\mathbf x,t)=\sum_{\lambda=1}^{2}
\int\frac{d^3k}{(2\pi)^3}\,
\boldsymbol{\epsilon}_{\lambda}(\mathbf k)\,
a_{\lambda}(\mathbf k)\cos\left(\mathbf k\cdot\mathbf x-\omega t
+\phi_{\lambda}(\mathbf k)\right),
\label{eq:RandomExpansion}
\end{equation}
where
\[\omega=c|\mathbf k|,\]
and
\[\boldsymbol{\epsilon}_{\lambda}\cdot\mathbf k=0\]
are transverse polarization vectors.
The amplitudes satisfy a prescribed spectral distribution, whereas the
phases
\[\phi_{\lambda}(\mathbf k)\]
are independent random variables uniformly distributed over
\[0\le\phi<2\pi.\]
This random-phase approximation is the central statistical assumption of
SED.
%%%%%%%%%%%%%%%%%%%%%%%%%%%%%%%%%%%%%%%%%%%%%%%%%%%%%%%%%%%%%%%%%%%%%%%%%%%%%%%%
\subsection{Random Variables}
%%%%%%%%%%%%%%%%%%%%%%%%%%%%%%%%%%%%%%%%%%%%%%%%%%%%%%%%%%%%%%%%%%%%%%%%%%%%%%%%
Ensemble averages are taken over the random phases.
Immediately,
\[\langle\cos(\theta+\phi)\rangle=0,\]
which implies
\begin{equation}
\boxed{\langle\mathbf E\rangle=\langle\mathbf B\rangle=0.}
\end{equation}
The vacuum therefore contains no preferred electric or magnetic field.
However,
quadratic averages do not vanish.
For example,
\[\langle E_iE_j\rangle\neq0,\]
indicating that the vacuum possesses finite electromagnetic
fluctuations despite zero mean fields.
This feature is completely analogous to ordinary Brownian motion, where
\[\langle dW\rangle=0,\qquad\langle(dW)^2\rangle=dt.\]
Thus,
the physically significant quantities are not the fields themselves but
their correlation functions.
%%%%%%%%%%%%%%%%%%%%%%%%%%%%%%%%%%%%%%%%%%%%%%%%%%%%%%%%%%%%%%%%%%%%%%%%%%%%%%%%
\subsection{Two-Point Correlation Functions}
%%%%%%%%%%%%%%%%%%%%%%%%%%%%%%%%%%%%%%%%%%%%%%%%%%%%%%%%%%%%%%%%%%%%%%%%%%%%%%%%
The statistical properties of the zero-point field are completely
specified by the two-point correlation functions,
\begin{equation}
C_{ij}(x,x')=\left<E_i(x)E_j(x')\right>.
\end{equation}
For a homogeneous and isotropic vacuum,
the correlation function depends only upon the coordinate difference,
\[x-x'.\]
Translation invariance therefore implies
\[C_{ij}(x,x')=C_{ij}(x-x').\]
Similarly,
rotational invariance restricts the tensor structure to combinations of
\[\delta_{ij},\qquad k_ik_j.\]
Lorentz invariance places even stronger constraints and ultimately fixes
the spectral dependence of the vacuum field.
%%%%%%%%%%%%%%%%%%%%%%%%%%%%%%%%%%%%%%%%%%%%%%%%%%%%%%%%%%%%%%%%%%%%%%%%%%%%%%%%
\subsection{Energy Density}
%%%%%%%%%%%%%%%%%%%%%%%%%%%%%%%%%%%%%%%%%%%%%%%%%%%%%%%%%%%%%%%%%%%%%%%%%%%%%%%%
The electromagnetic energy density is
\begin{equation}
u=\frac1{8\pi}(E^2+B^2)
\end{equation}
(in Gaussian units).
Taking the ensemble average gives
\[\langle u\rangle=\frac1{8\pi}\left<E^2+B^2\right>.\]
Using the plane-wave expansion yields
\begin{equation}
\langle u\rangle=\int_0^\infty u(\omega)\,d\omega,
\end{equation}
where
\[u(\omega)\]
is the spectral energy density.
Determining this function constitutes the central problem of SED.
%%%%%%%%%%%%%%%%%%%%%%%%%%%%%%%%%%%%%%%%%%%%%%%%%%%%%%%%%%%%%%%%%%%%%%%%%%%%%%%%
\subsection{Why the Spectrum Is Not Arbitrary}
%%%%%%%%%%%%%%%%%%%%%%%%%%%%%%%%%%%%%%%%%%%%%%%%%%%%%%%%%%%%%%%%%%%%%%%%%%%%%%%%
At first sight one might imagine choosing any spectral density.
This is impossible.
The vacuum must satisfy simultaneously
\begin{enumerate}
\item Maxwell's equations,
\item statistical homogeneity,
\item isotropy,
\item Lorentz invariance,
\item scale invariance (at zero temperature),
\item positivity of the energy density.
\end{enumerate}
These requirements are so restrictive that only one spectral form
survives (up to an overall multiplicative constant).
As first emphasized by Boyer,
the Lorentz-invariant zero-temperature spectrum is proportional to
\[\omega^3.\]
The overall proportionality constant is then identified empirically with
\[\frac{\hbar}{2\pi^2c^3},\]
giving
\begin{equation}
\boxed{u(\omega)=\frac{\hbar\omega^3}{2\pi^2c^3}.}
\label{eq:BoyerSpectrum}
\end{equation}
This result is among the most remarkable aspects of SED.
Planck's constant appears not because the field is quantized, but as the
measured intensity of a classical Lorentz-invariant stochastic vacuum.
%%%%%%%%%%%%%%%%%%%%%%%%%%%%%%%%%%%%%%%%%%%%%%%%%%%%%%%%%%%%%%%%%%%%%%%%%%%%%%%%
\subsection{The Origin of Planck's Constant}
%%%%%%%%%%%%%%%%%%%%%%%%%%%%%%%%%%%%%%%%%%%%%%%%%%%%%%%%%%%%%%%%%%%%%%%%%%%%%%%%
Equation (\ref{eq:BoyerSpectrum}) raises an immediate conceptual
question.
Why should a classical theory contain the quantum constant
$\hbar$?
SED answers this question differently from ordinary quantum mechanics.
Quantum mechanics assumes that $\hbar$ is a fundamental constant of
nature entering directly into the commutation relations,
\[[x,p]=i\hbar.\]
SED assumes instead that $\hbar$ merely determines the intensity of the
classical zero-point field.
From this perspective,
the appearance of Planck's constant is no more mysterious than the
appearance of the speed of light in Maxwell's equations.
Both are universal constants characterizing the vacuum.
%%%%%%%%%%%%%%%%%%%%%%%%%%%%%%%%%%%%%%%%%%%%%%%%%%%%%%%%%%%%%%%%%%%%%%%%%%%%%%%%
\subsection*{Critical Assessment}
%%%%%%%%%%%%%%%%%%%%%%%%%%%%%%%%%%%%%%%%%%%%%%%%%%%%%%%%%%%%%%%%%%%%%%%%%%%%%%%%
The mathematical construction of the zero-point field is internally
consistent and remarkably economical. Starting from classical Maxwell
theory, together with assumptions of homogeneity, isotropy and Lorentz
invariance, one is naturally led to a unique vacuum spectrum whose form
closely resembles that suggested by quantum theory.
Nevertheless, several important conceptual questions remain open.
The normalization of the spectrum cannot be derived within classical
electrodynamics alone and must be fixed empirically through the observed
value of Planck's constant.
Moreover, the total vacuum energy,
\[\int_0^\infty\frac{\hbar\omega^3}{2\pi^2c^3}\,d\omega,\]
diverges quartically. In practice this divergence is handled by
considering energy differences, introducing physical cutoffs, or
restricting attention to observable correlation functions, much as in
quantum field theory. The physical interpretation of the absolute vacuum
energy remains one of the major unresolved issues shared by both SED and
conventional quantum electrodynamics.
From the perspective of the present monograph, the most important lesson
is that stochastic electrodynamics introduces randomness through the
vacuum field rather than through particle trajectories. This alternative
origin of stochasticity will provide an instructive point of comparison
with Nelson's diffusion theory and, later, with the stochastic
Hamiltonian formulation developed in the SHP framework.
%%%%%%%%%%%%%%%%%%%%%%%%%%%%%%%%%%%%%%%%%%%%%%%%%%%%%%%%%%%%%%%%%%%%%%%%%%%%%%%%
\section{Derivation of the Lorentz-Invariant Zero-Point Spectrum}
\label{sec:LorentzInvariantSpectrum}
%%%%%%%%%%%%%%%%%%%%%%%%%%%%%%%%%%%%%%%%%%%%%%%%%%%%%%%%%%%%%%%%%%%%%%%%%%%%%%%%
The cornerstone of stochastic electrodynamics is the assumption that the
vacuum contains a homogeneous, isotropic stochastic electromagnetic
field. Unlike ordinary thermal radiation, this field must exist even at
zero temperature and must not single out any preferred inertial frame.
These requirements are sufficiently restrictive that the spectral energy
density is determined almost uniquely. In this section we derive the
Lorentz-invariant spectrum following the physical arguments originally
developed by Boyer and later refined by several authors.
%%%%%%%%%%%%%%%%%%%%%%%%%%%%%%%%%%%%%%%%%%%%%%%%%%%%%%%%%%%%%%%%%%%%%%%%%%%%%%%%
\subsection{General Spectral Representation}
%%%%%%%%%%%%%%%%%%%%%%%%%%%%%%%%%%%%%%%%%%%%%%%%%%%%%%%%%%%%%%%%%%%%%%%%%%%%%%%%
Let
\[u(\omega)\,d\omega\]
denote the electromagnetic energy per unit volume contained in the
frequency interval
\[(\omega,\omega+d\omega).\]
The total vacuum energy density is therefore
\begin{equation}
u=\int_0^\infty u(\omega)\,d\omega.
\end{equation}
At this stage the function
\[\nu(\omega)\]
is completely arbitrary.
Our task is to determine its form using only general physical
principles.
%%%%%%%%%%%%%%%%%%%%%%%%%%%%%%%%%%%%%%%%%%%%%%%%%%%%%%%%%%%%%%%%%%%%%%%%%%%%%%%%
\subsection{Dimensional Analysis}
%%%%%%%%%%%%%%%%%%%%%%%%%%%%%%%%%%%%%%%%%%%%%%%%%%%%%%%%%%%%%%%%%%%%%%%%%%%%%%%%
The dimensions of the spectral density are
\[[\nu]=\frac{\mathrm{Energy}}{\mathrm{Volume}\times\mathrm{Frequency}}.\]
Using
\[[\omega]=T^{-1},\qquad[c]=LT^{-1},\]
one immediately finds that the only possible Lorentz-covariant
combination having the correct dimensions is
\begin{equation}
\nu(\omega)=C\,\omega^3,
\end{equation}
where
\[C\]
is a universal constant having dimensions of action divided by
\(c^3\).
Dimensional analysis alone therefore already suggests a cubic spectrum.
However, dimensional analysis cannot determine the constant
\(C\), nor can it prove Lorentz invariance.
A more detailed argument is required.
%%%%%%%%%%%%%%%%%%%%%%%%%%%%%%%%%%%%%%%%%%%%%%%%%%%%%%%%%%%%%%%%%%%%%%%%%%%%%%%%
\subsection{Lorentz Transformation of Plane Waves}
%%%%%%%%%%%%%%%%%%%%%%%%%%%%%%%%%%%%%%%%%%%%%%%%%%%%%%%%%%%%%%%%%%%%%%%%%%%%%%%%
Consider a monochromatic plane wave with four-wavevector
\[k^\mu=\left(\frac{\omega}{c},\mathbf k\right),\qquad k^\mu k_\mu=0.\]
Under a Lorentz transformation,
\[k^\mu\longrightarrow k'^\mu=\Lambda^\mu{}_\nu k^\nu .\]
The observed frequency changes according to the Doppler formula,
\begin{equation}
\omega'=\gamma(1-\beta\cos\theta)\,\omega .
\label{eq:Doppler}
\end{equation}
Consequently,
both the frequency and the density of states transform under boosts.
A Lorentz-invariant vacuum requires that these two effects compensate
each other exactly.
%%%%%%%%%%%%%%%%%%%%%%%%%%%%%%%%%%%%%%%%%%%%%%%%%%%%%%%%%%%%%%%%%%%%%%%%%%%%%%%%
\subsection{Density of Electromagnetic Modes}
%%%%%%%%%%%%%%%%%%%%%%%%%%%%%%%%%%%%%%%%%%%%%%%%%%%%%%%%%%%%%%%%%%%%%%%%%%%%%%%%
In a volume
\[V,\]
the number of electromagnetic modes contained in the interval
\[(\omega,\omega+d\omega)\]
is
\begin{equation}
g(\omega)d\omega=\frac{\omega^2}{\pi^2c^3}\,d\omega,
\label{eq:ModeDensity}
\end{equation}
including both polarization states.
The spectral energy density may therefore be written as
\begin{equation}
\nu(\omega)=g(\omega)\,\varepsilon(\omega),
\end{equation}
where
\[\varepsilon(\omega)\]
is the average energy carried by each mode.
%%%%%%%%%%%%%%%%%%%%%%%%%%%%%%%%%%%%%%%%%%%%%%%%%%%%%%%%%%%%%%%%%%%%%%%%%%%%%%%%
\subsection{Scale Invariance of the Vacuum}
%%%%%%%%%%%%%%%%%%%%%%%%%%%%%%%%%%%%%%%%%%%%%%%%%%%%%%%%%%%%%%%%%%%%%%%%%%%%%%%%
At zero temperature the vacuum contains no intrinsic frequency scale.
The spectrum must therefore remain invariant under the simultaneous
transformation
\[\omega\rightarrow\lambda\omega .\]
Suppose
\[\varepsilon(\omega)\propto\omega^n .\]
Then
\[\nu(\omega)\propto\omega^{n+2}.\]
Lorentz invariance of the vacuum stress-energy tensor requires
\[n=1,\]
leading immediately to
\begin{equation}
\boxed{\nu(\omega)\propto\omega^3.}
\end{equation}
Thus the cubic spectrum is not an arbitrary assumption.
It follows directly from Lorentz symmetry together with the absence of
any preferred frequency scale.
%%%%%%%%%%%%%%%%%%%%%%%%%%%%%%%%%%%%%%%%%%%%%%%%%%%%%%%%%%%%%%%%%%%%%%%%%%%%%%%%
\subsection{Normalization}
%%%%%%%%%%%%%%%%%%%%%%%%%%%%%%%%%%%%%%%%%%%%%%%%%%%%%%%%%%%%%%%%%%%%%%%%%%%%%%%%
The proportionality constant cannot be determined by classical
electrodynamics.
Experiment fixes it to be
\begin{equation}
\boxed{\nu(\omega)=\frac{\hbar\omega^3}{2\pi^2c^3}.}
\label{eq:FinalSpectrum}
\end{equation}
Equivalently,
the average energy per electromagnetic normal mode becomes
\begin{equation}
\boxed{\varepsilon(\omega)=\frac12\hbar\omega.}
\label{eq:HalfQuantum}
\end{equation}
This expression is numerically identical to the zero-point energy of a
quantum harmonic oscillator.
In SED, however, its interpretation is entirely different.
It is regarded as the energy carried by each classical stochastic mode
of the electromagnetic vacuum.
%%%%%%%%%%%%%%%%%%%%%%%%%%%%%%%%%%%%%%%%%%%%%%%%%%%%%%%%%%%%%%%%%%%%%%%%%%%%%%%%
\subsection{Vacuum Correlation Functions}
%%%%%%%%%%%%%%%%%%%%%%%%%%%%%%%%%%%%%%%%%%%%%%%%%%%%%%%%%%%%%%%%%%%%%%%%%%%%%%%%
The spectral density determines all vacuum correlation functions.
For example,
\begin{equation}
\langle E_i(\mathbf x,t)E_j(\mathbf x',t')\rangle=
\int d^3k\,P_{ij}(\mathbf k)\,\nu(\omega)\cos
\left[\mathbf k\cdot(\mathbf x-\mathbf x')-\omega(t-t')\right],
\end{equation}
where
\[P_{ij}=\delta_{ij}-\frac{k_ik_j}{k^2}\]
is the transverse projection operator.
These correlation functions provide the stochastic forces responsible
for all particle dynamics in stochastic electrodynamics.
%%%%%%%%%%%%%%%%%%%%%%%%%%%%%%%%%%%%%%%%%%%%%%%%%%%%%%%%%%%%%%%%%%%%%%%%%%%%%%%%
\subsection{Comparison with Quantum Electrodynamics}
%%%%%%%%%%%%%%%%%%%%%%%%%%%%%%%%%%%%%%%%%%%%%%%%%%%%%%%%%%%%%%%%%%%%%%%%%%%%%%%%
The formal similarity between Eqs.~(\ref{eq:FinalSpectrum}) and
(\ref{eq:HalfQuantum}) and the corresponding expressions of quantum
electrodynamics is striking.
Nevertheless, the physical interpretations differ profoundly.
In quantum electrodynamics,
\[\frac12\hbar\omega\]
represents the ground-state energy of a quantized harmonic oscillator.
In stochastic electrodynamics,
the same quantity represents the average energy contained in a classical
random electromagnetic mode.
Thus identical mathematical expressions arise from two entirely
different ontological assumptions.
%%%%%%%%%%%%%%%%%%%%%%%%%%%%%%%%%%%%%%%%%%%%%%%%%%%%%%%%%%%%%%%%%%%%%%%%%%%%%%%%
\subsection*{Critical Assessment}
%%%%%%%%%%%%%%%%%%%%%%%%%%%%%%%%%%%%%%%%%%%%%%%%%%%%%%%%%%%%%%%%%%%%%%%%%%%%%%%%
The derivation above illustrates both the strength and the weakness of
stochastic electrodynamics.
On the one hand, the cubic spectrum follows naturally from very general
symmetry arguments and reproduces one of the most characteristic
numerical features of quantum vacuum fluctuations.
On the other hand, the overall normalization remains empirical. The
appearance of Planck's constant is not explained but assumed through the
measured intensity of the vacuum field. Consequently, SED shifts the
fundamental question from the quantization of matter to the origin of
the stochastic vacuum itself.
Nevertheless, the existence of a unique Lorentz-invariant classical
vacuum spectrum remains one of the most elegant and physically
compelling aspects of the theory.
%%%%%%%%%%%%%%%%%%%%%%%%%%%%%%%%%%%%%%%%%%%%%%%%%%%%%%%%%%%%%%%%%%%%%%%%%%%%%%%%
\section{Radiation Reaction and the Fluctuation--Dissipation Balance}
\label{sec:RadiationReaction}
%%%%%%%%%%%%%%%%%%%%%%%%%%%%%%%%%%%%%%%%%%%%%%%%%%%%%%%%%%%%%%%%%%%%%%%%%%%%%%%%
The existence of a stochastic electromagnetic vacuum is not, by itself,
sufficient to explain quantum phenomena. One must also understand how a
charged particle responds to this fluctuating field.
The dynamics of a charged particle differs fundamentally from ordinary
Brownian motion. Besides the random force exerted by the vacuum field,
an accelerated charge continuously emits electromagnetic radiation.
Energy is therefore lost through radiation reaction.
The motion is governed by the competition between two opposing effects:
\[\boxed{\text{Vacuum fluctuations}\Longleftrightarrow\text{Radiation reaction}.}\]
The former inject energy into the particle, whereas the latter removes
it. A stationary state becomes possible only if these two processes
balance each other.
%%%%%%%%%%%%%%%%%%%%%%%%%%%%%%%%%%%%%%%%%%%%%%%%%%%%%%%%%%%%%%%%%%%%%%%%%%%%%%%%
\subsection{Lorentz Force}
%%%%%%%%%%%%%%%%%%%%%%%%%%%%%%%%%%%%%%%%%%%%%%%%%%%%%%%%%%%%%%%%%%%%%%%%%%%%%%%%
The relativistic equation of motion is
\begin{equation}
m\frac{du^\mu}{d\tau}=eF^{\mu\nu}u_\nu ,
\label{eq:LorentzForce}
\end{equation}
where
\[F^{\mu\nu}=F^{\mu\nu}_{\rm ext}+F^{\mu\nu}_{\rm ZP}.\]
The second term denotes the stochastic zero-point field.
Equation (\ref{eq:LorentzForce}) is entirely classical.
Randomness enters only through the stochastic electromagnetic tensor.
%%%%%%%%%%%%%%%%%%%%%%%%%%%%%%%%%%%%%%%%%%%%%%%%%%%%%%%%%%%%%%%%%%%%%%%%%%%%%%%%
\subsection{Radiation Reaction}
%%%%%%%%%%%%%%%%%%%%%%%%%%%%%%%%%%%%%%%%%%%%%%%%%%%%%%%%%%%%%%%%%%%%%%%%%%%%%%%%
Accelerated charges radiate electromagnetic energy.
Conservation of energy therefore requires a self-force acting on the
particle.
In the nonrelativistic approximation,
\begin{equation}
m\ddot x=F_{\rm ext}+m\tau_0\dddot x ,
\label{eq:AbrahamLorentz}
\end{equation}
where
\begin{equation}
\tau_0=\frac{2e^2}{3mc^3}.
\end{equation}
For an electron,
\[\tau_0\simeq 6.3\times10^{-24}\ {\rm s},\]
an extremely short time scale.
The third-order derivative reflects the fact that radiation reaction
depends upon the rate of change of acceleration.
%%%%%%%%%%%%%%%%%%%%%%%%%%%%%%%%%%%%%%%%%%%%%%%%%%%%%%%%%%%%%%%%%%%%%%%%%%%%%%%%
\subsection{Relativistic Form}
%%%%%%%%%%%%%%%%%%%%%%%%%%%%%%%%%%%%%%%%%%%%%%%%%%%%%%%%%%%%%%%%%%%%%%%%%%%%%%%%
The fully covariant Abraham--Lorentz--Dirac equation is
\begin{equation}
m\frac{du^\mu}{d\tau}=eF^{\mu\nu}u_\nu+\frac{2e^2}{3c^3}
\left(\frac{d^2u^\mu}{d\tau^2}+u^\mu\frac{du_\nu}{d\tau}\frac{du^\nu}{d\tau}\right).
\label{eq:ALD}
\end{equation}
The second term is the radiation-reaction force.
Although derived entirely within classical electrodynamics, it has
generated considerable controversy because of its unusual mathematical
properties.
%%%%%%%%%%%%%%%%%%%%%%%%%%%%%%%%%%%%%%%%%%%%%%%%%%%%%%%%%%%%%%%%%%%%%%%%%%%%%%%%
\subsection{Runaway Solutions}
%%%%%%%%%%%%%%%%%%%%%%%%%%%%%%%%%%%%%%%%%%%%%%%%%%%%%%%%%%%%%%%%%%%%%%%%%%%%%%%%
Equation (\ref{eq:ALD}) possesses exponentially growing solutions,
\[a(t)\propto e^{t/\tau_0},\]
even when no external force acts.
These runaway solutions are clearly unphysical.
They arise because the equation is third order in time.
%%%%%%%%%%%%%%%%%%%%%%%%%%%%%%%%%%%%%%%%%%%%%%%%%%%%%%%%%%%%%%%%%%%%%%%%%%%%%%%%
\subsection{Preacceleration}
%%%%%%%%%%%%%%%%%%%%%%%%%%%%%%%%%%%%%%%%%%%%%%%%%%%%%%%%%%%%%%%%%%%%%%%%%%%%%%%%
A second difficulty is preacceleration.
Some solutions begin to accelerate before the external force is applied,
\[a(t)\neq0,\qquad t<t_{\rm force}.\]
This apparent violation of causality has stimulated an extensive
literature and remains one of the most discussed aspects of classical
electrodynamics.
%%%%%%%%%%%%%%%%%%%%%%%%%%%%%%%%%%%%%%%%%%%%%%%%%%%%%%%%%%%%%%%%%%%%%%%%%%%%%%%%
\subsection{Landau--Lifshitz Approximation}
%%%%%%%%%%%%%%%%%%%%%%%%%%%%%%%%%%%%%%%%%%%%%%%%%%%%%%%%%%%%%%%%%%%%%%%%%%%%%%%%
For slowly varying external fields one may eliminate these pathological
solutions perturbatively.
Replacing the acceleration appearing inside the radiation-reaction term
by the Lorentz force yields the Landau--Lifshitz equation,
\begin{equation}
m\frac{du^\mu}{d\tau}=eF^{\mu\nu}u_\nu+\frac{2e^3}{3mc^3}\partial_\alpha 
    F^{\mu\nu}u^\alpha u_\nu+\cdots .
\end{equation}
This equation is second order and avoids both runaway solutions and
preacceleration within its domain of validity.
For most practical applications it has become the preferred classical
equation of motion.
%%%%%%%%%%%%%%%%%%%%%%%%%%%%%%%%%%%%%%%%%%%%%%%%%%%%%%%%%%%%%%%%%%%%%%%%%%%%%%%%
\subsection{Fluctuation--Dissipation Balance}
%%%%%%%%%%%%%%%%%%%%%%%%%%%%%%%%%%%%%%%%%%%%%%%%%%%%%%%%%%%%%%%%%%%%%%%%%%%%%%%%
The stochastic vacuum continuously transfers energy to the charged
particle,
\[\left<P_{\rm in}\right>>0.\]
Radiation reaction removes energy,
\[\left<P_{\rm out}\right>>0.\]
A stationary equilibrium requires
\begin{equation}
\boxed{\left<P_{\rm in}\right>=\left<P_{\rm out}\right>.}
\label{eq:FDT}
\end{equation}
This relation is the SED analogue of the fluctuation--dissipation
theorem familiar from statistical mechanics.
Unlike ordinary Brownian motion, however, both fluctuation and
dissipation originate from classical electrodynamics itself.
%%%%%%%%%%%%%%%%%%%%%%%%%%%%%%%%%%%%%%%%%%%%%%%%%%%%%%%%%%%%%%%%%%%%%%%%%%%%%%%%
\subsection{Relation to Brownian Motion}
%%%%%%%%%%%%%%%%%%%%%%%%%%%%%%%%%%%%%%%%%%%%%%%%%%%%%%%%%%%%%%%%%%%%%%%%%%%%%%%%
The analogy with Langevin dynamics is illuminating.
For an ordinary Brownian particle,
\[m\dot v=-\gamma v+\xi(t),\]
where
\[\xi(t)\]
is the stochastic force and
\[\gamma\]
represents viscous damping.
In stochastic electrodynamics,
\[\xi(t)\longrightarrow eE_{\rm ZP}(t),\]
while the damping term is replaced by radiation reaction.
Thus the entire mathematical structure of Brownian motion reappears in a
purely electromagnetic context.
%%%%%%%%%%%%%%%%%%%%%%%%%%%%%%%%%%%%%%%%%%%%%%%%%%%%%%%%%%%%%%%%%%%%%%%%%%%%%%%%
\subsection{Conceptual Significance}
%%%%%%%%%%%%%%%%%%%%%%%%%%%%%%%%%%%%%%%%%%%%%%%%%%%%%%%%%%%%%%%%%%%%%%%%%%%%%%%%
The fluctuation--dissipation balance occupies a central position in
stochastic electrodynamics.
Without radiation reaction the stochastic vacuum would continuously
increase the particle energy.
Without vacuum fluctuations the particle would radiate away all of its
energy and eventually come to rest.
Only the simultaneous presence of both mechanisms permits stable
stationary states.
This balance is responsible for many of the characteristic predictions
of SED, including the ground-state energy of the harmonic oscillator.
%%%%%%%%%%%%%%%%%%%%%%%%%%%%%%%%%%%%%%%%%%%%%%%%%%%%%%%%%%%%%%%%%%%%%%%%%%%%%%%%
\subsection*{Critical Assessment}
%%%%%%%%%%%%%%%%%%%%%%%%%%%%%%%%%%%%%%%%%%%%%%%%%%%%%%%%%%%%%%%%%%%%%%%%%%%%%%%%
The interplay between vacuum fluctuations and radiation reaction is one
of the strongest aspects of stochastic electrodynamics because both
ingredients arise naturally within classical electrodynamics.
At the same time, radiation reaction remains one of the oldest
conceptual problems in theoretical physics. The Abraham--Lorentz--Dirac
equation exhibits runaway solutions and preacceleration, while the
Landau--Lifshitz equation is an approximation rather than an exact
theory.
Consequently, the foundations of SED inherit the unresolved difficulties
of classical radiation reaction. Whether these problems reflect a
limitation of classical electrodynamics itself or merely indicate the
need for a more complete microscopic description remains an open
question.
From the perspective of this monograph, an important lesson emerges:
the stochastic force and the dissipative self-force are not independent
ingredients but two complementary aspects of the interaction between a
charged particle and its electromagnetic environment. This observation
will later prove useful when constructing stochastic Hamiltonian
dynamics within the SHP framework.
%%%%%%%%%%%%%%%%%%%%%%%%%%%%%%%%%%%%%%%%%%%%%%%%%%%%%%%%%%%%%%%%%%%%%%%%%%%%%%%%
\section{The Harmonic Oscillator in Stochastic Electrodynamics}
\label{sec:SEDOscillator}
%%%%%%%%%%%%%%%%%%%%%%%%%%%%%%%%%%%%%%%%%%%%%%%%%%%%%%%%%%%%%%%%%%%%%%%%%%%%%%%%
The one-dimensional harmonic oscillator occupies a central position in
both classical and quantum physics. Nearly every physical system can be
approximated by a harmonic oscillator near a stable equilibrium.
Moreover, the electromagnetic field itself is equivalent to an infinite
collection of independent harmonic oscillators. Consequently, any
candidate stochastic foundation of quantum mechanics must reproduce the
properties of the harmonic oscillator before more complicated systems
can be considered.
For stochastic electrodynamics, the harmonic oscillator provides the
first decisive test of the theory.
%%%%%%%%%%%%%%%%%%%%%%%%%%%%%%%%%%%%%%%%%%%%%%%%%%%%%%%%%%%%%%%%%%%%%%%%%%%%%%%%
\subsection{Equation of Motion}
%%%%%%%%%%%%%%%%%%%%%%%%%%%%%%%%%%%%%%%%%%%%%%%%%%%%%%%%%%%%%%%%%%%%%%%%%%%%%%%%
Consider a particle of mass $m$ and charge $e$ bound in the harmonic
potential
\begin{equation}
V(x)=\frac12 m\omega_0^2x^2.
\end{equation}
Including both the stochastic vacuum field and radiation reaction, the
equation of motion becomes
\begin{equation}
m\ddot x+m\omega_0^2x=eE_{\rm ZP}(t)+m\tau_0\dddot x,
\label{eq:SEDOscillator}
\end{equation}
where
\[\tau_0=\frac{2e^2}{3mc^3}.\]
Equation (\ref{eq:SEDOscillator}) is simply the Abraham--Lorentz
equation for a harmonically bound particle.
The stochastic force originates entirely from the classical zero-point
field.
%%%%%%%%%%%%%%%%%%%%%%%%%%%%%%%%%%%%%%%%%%%%%%%%%%%%%%%%%%%%%%%%%%%%%%%%%%%%%%%%
\subsection{Frequency-Domain Solution}
%%%%%%%%%%%%%%%%%%%%%%%%%%%%%%%%%%%%%%%%%%%%%%%%%%%%%%%%%%%%%%%%%%%%%%%%%%%%%%%%
Because the stochastic force contains all frequencies, it is convenient
to Fourier transform Eq.~(\ref{eq:SEDOscillator}).
Writing
\[x(t)=\int_{-\infty}^{\infty}\tilde x(\omega)e^{-i\omega t}\,d\omega,\]
one obtains
\begin{equation}
\left(\omega_0^2-\omega^2-i\tau_0\omega^3\right)\tilde x(\omega)=
\frac{e}{m}\tilde E(\omega).
\end{equation}
Hence,
\begin{equation}
\boxed{\tilde x(\omega)=
\frac{e/m}{\omega_0^2-\omega^2-i\tau_0\omega^3}\,\tilde E(\omega).}
\label{eq:ResponseFunction}
\end{equation}
The denominator is the linear response function of the oscillator.
The imaginary part represents energy dissipation through radiation
reaction.
%%%%%%%%%%%%%%%%%%%%%%%%%%%%%%%%%%%%%%%%%%%%%%%%%%%%%%%%%%%%%%%%%%%%%%%%%%%%%%%%
\subsection{Power Spectrum}
%%%%%%%%%%%%%%%%%%%%%%%%%%%%%%%%%%%%%%%%%%%%%%%%%%%%%%%%%%%%%%%%%%%%%%%%%%%%%%%%
The displacement spectrum follows immediately,
\begin{equation}
S_x(\omega)=\left|\frac{e/m}{\omega_0^2-\omega^2-i\tau_0\omega^3}\right|^2 
    S_E(\omega),
\label{eq:PowerSpectrum}
\end{equation}
where
\[S_E(\omega)\]
is the spectral density of the zero-point field.
Using the Lorentz-invariant spectrum derived in the previous section,
\begin{equation}
S_E(\omega)\propto\hbar\omega^3.
\end{equation}
Thus every normal mode of the vacuum contributes to the oscillator
motion, although the resonance at
\[\omega=\omega_0\]
dominates because of the response function.
%%%%%%%%%%%%%%%%%%%%%%%%%%%%%%%%%%%%%%%%%%%%%%%%%%%%%%%%%%%%%%%%%%%%%%%%%%%%%%%%
\subsection{Mean-Square Displacement}
%%%%%%%%%%%%%%%%%%%%%%%%%%%%%%%%%%%%%%%%%%%%%%%%%%%%%%%%%%%%%%%%%%%%%%%%%%%%%%%%
The variance of the position is
\begin{equation}
\langle x^2\rangle=\int_0^\infty S_x(\omega)\,d\omega.
\label{eq:XVarianceIntegral}
\end{equation}
The integral is sharply peaked near the resonance frequency.
Since
\[\tau_0\omega_0\ll1,\]
one may evaluate the integral using the narrow-resonance approximation.
Expanding about
\[\omega=\omega_0,\]
and retaining the leading contribution yields
\begin{equation}
\boxed{\langle x^2\rangle=\frac{\hbar}{2m\omega_0}.}
\label{eq:GroundVariance}
\end{equation}
Remarkably, this is exactly the quantum-mechanical ground-state
expectation value.
%%%%%%%%%%%%%%%%%%%%%%%%%%%%%%%%%%%%%%%%%%%%%%%%%%%%%%%%%%%%%%%%%%%%%%%%%%%%%%%%
\subsection{Momentum Fluctuations}
%%%%%%%%%%%%%%%%%%%%%%%%%%%%%%%%%%%%%%%%%%%%%%%%%%%%%%%%%%%%%%%%%%%%%%%%%%%%%%%%
Since
\[p=m\dot x,\]
the momentum variance becomes
\begin{equation}
\langle p^2\rangle=m^2\int_0^\infty\omega^2 S_x(\omega)\,d\omega.
\end{equation}
The same approximation gives
\begin{equation}
\boxed{\langle p^2\rangle=\frac12 m\hbar\omega_0.}
\label{eq:MomentumVariance}
\end{equation}
Consequently,
\begin{equation}
\Delta x \Delta p=\frac{\hbar}{2},
\end{equation}
which coincides with the minimum uncertainty relation of quantum
mechanics.
Within SED this relation emerges from stochastic dynamics rather than
from operator commutation relations.
%%%%%%%%%%%%%%%%%%%%%%%%%%%%%%%%%%%%%%%%%%%%%%%%%%%%%%%%%%%%%%%%%%%%%%%%%%%%%%%%
\subsection{Mean Energy}
%%%%%%%%%%%%%%%%%%%%%%%%%%%%%%%%%%%%%%%%%%%%%%%%%%%%%%%%%%%%%%%%%%%%%%%%%%%%%%%%
The average oscillator energy is
\begin{equation}
\langle E\rangle=\frac{\langle p^2\rangle}{2m}+\frac12 m\omega_0^2
    \langle x^2\rangle .
\end{equation}
Using Eqs.~(\ref{eq:GroundVariance}) and
(\ref{eq:MomentumVariance}) gives
\begin{equation}
\boxed{\langle E\rangle=\frac12\hbar\omega_0.}
\label{eq:GroundEnergy}
\end{equation}
Thus the classical stochastic oscillator possesses precisely the same
average ground-state energy as the quantum oscillator.
%%%%%%%%%%%%%%%%%%%%%%%%%%%%%%%%%%%%%%%%%%%%%%%%%%%%%%%%%%%%%%%%%%%%%%%%%%%%%%%%
\subsection{Probability Distribution}
%%%%%%%%%%%%%%%%%%%%%%%%%%%%%%%%%%%%%%%%%%%%%%%%%%%%%%%%%%%%%%%%%%%%%%%%%%%%%%%%
The stochastic force is Gaussian.
Because the oscillator responds linearly, the stationary probability
distribution of the displacement is also Gaussian,
\begin{equation}
P(x)=\frac1{\sqrt{2\pi\langle x^2\rangle}}
\exp\left(-\frac{x^2}{2\langle x^2\rangle}\right).
\end{equation}
Substituting Eq.~(\ref{eq:GroundVariance}) yields
\begin{equation}
\boxed{P(x)=\left(\frac{m\omega_0}{\pi\hbar}\right)^{1/2}
\exp\left(-\frac{m\omega_0x^2}{\hbar}\right),}
\label{eq:GaussianSED}
\end{equation}
which is identical to the probability density obtained from the quantum
ground-state wave function,
\[|\psi_0(x)|^2.\]
%%%%%%%%%%%%%%%%%%%%%%%%%%%%%%%%%%%%%%%%%%%%%%%%%%%%%%%%%%%%%%%%%%%%%%%%%%%%%%%%
\subsection{Correlation Functions}
%%%%%%%%%%%%%%%%%%%%%%%%%%%%%%%%%%%%%%%%%%%%%%%%%%%%%%%%%%%%%%%%%%%%%%%%%%%%%%%%
The stationary two-time correlation function is
\begin{equation}
C(\tau)=\langle x(t)x(t+\tau)\rangle .
\end{equation}
Using the Wiener--Khinchin theorem,
\begin{equation}
C(\tau)=\int_0^\infty S_x(\omega)\cos(\omega\tau)\,d\omega.
\end{equation}
Near resonance,
\begin{equation}
C(\tau)\simeq\frac{\hbar}{2m\omega_0}\cos(\omega_0\tau)e^{-\Gamma|\tau|},
\end{equation}
where
\[\Gamma=\frac12\tau_0\omega_0^2\]
is the radiative damping rate.
Thus the stochastic oscillator possesses a finite coherence time,
determined by radiation reaction.
%%%%%%%%%%%%%%%%%%%%%%%%%%%%%%%%%%%%%%%%%%%%%%%%%%%%%%%%%%%%%%%%%%%%%%%%%%%%%%%%
\subsection{Physical Interpretation}
%%%%%%%%%%%%%%%%%%%%%%%%%%%%%%%%%%%%%%%%%%%%%%%%%%%%%%%%%%%%%%%%%%%%%%%%%%%%%%%%
The agreement with the quantum ground state is one of the principal
achievements of stochastic electrodynamics.
No quantization has been imposed.
The oscillator obeys a purely classical equation of motion driven by a
classical stochastic electromagnetic field.
The quantities usually interpreted as consequences of canonical
commutation relations arise here from the statistical equilibrium
between vacuum fluctuations and radiative damping.
From the SED perspective, the ground-state energy is not the zero-point
energy of a quantized oscillator but the equilibrium energy of a
classical oscillator immersed in the universal zero-point radiation
field.
%%%%%%%%%%%%%%%%%%%%%%%%%%%%%%%%%%%%%%%%%%%%%%%%%%%%%%%%%%%%%%%%%%%%%%%%%%%%%%%%
\subsection*{Critical Assessment}
%%%%%%%%%%%%%%%%%%%%%%%%%%%%%%%%%%%%%%%%%%%%%%%%%%%%%%%%%%%%%%%%%%%%%%%%%%%%%%%%
The harmonic oscillator constitutes the strongest single success of
stochastic electrodynamics. The agreement with the quantum-mechanical
ground state extends beyond the average energy to include the position
variance, momentum variance, Gaussian probability distribution, and
minimum uncertainty product.
Nevertheless, several qualifications are necessary.
First, the derivation relies on the assumed Lorentz-invariant
zero-point spectrum introduced earlier; the normalization of that
spectrum remains an empirical input. Second, the calculation describes
the stationary ground state of a linear system. The extension to excited
states and strongly nonlinear systems is considerably more difficult.
Finally, the analysis is based on classical radiation reaction, whose
foundations remain controversial because of the Abraham–Lorentz–Dirac
pathologies discussed in the previous section.
Despite these limitations, the harmonic oscillator remains compelling
evidence that a classical stochastic field can reproduce certain
characteristic features traditionally regarded as uniquely quantum.
Whether this agreement reflects a deep physical mechanism or a
remarkable mathematical coincidence is one of the central questions
addressed by the subsequent development of stochastic electrodynamics.
%%%%%%%%%%%%%%%%%%%%%%%%%%%%%%%%%%%%%%%%%%%%%%%%%%%%%%%%%%%%%%%%%%%%%%%%%%%%%%%%
\section{The Hydrogen Atom}
\label{sec:SEDHydrogen}
%%%%%%%%%%%%%%%%%%%%%%%%%%%%%%%%%%%%%%%%%%%%%%%%%%%%%%%%%%%%%%%%%%%%%%%%%%%%%%%%
The hydrogen atom has always occupied a central position in the
development of atomic physics. Historically, one of the principal
motivations for quantum mechanics was the apparent inability of
classical electrodynamics to explain the stability of atomic matter.
For this reason the hydrogen atom provides the most stringent test of
stochastic electrodynamics. Unlike the harmonic oscillator, whose
equations are linear, the Coulomb problem is nonlinear, non-compact, and
contains trajectories with arbitrarily high eccentricity. Moreover, the
electron continuously accelerates in the Coulomb field and therefore
radiates according to classical electrodynamics.
The central question addressed by stochastic electrodynamics may be
stated as follows:
\begin{quote}
Can the interaction between a classical electron and the stochastic
electromagnetic vacuum prevent the classical radiative collapse of the
hydrogen atom?
\end{quote}
Unlike the harmonic oscillator, this question has not yet received a
universally accepted answer.
%%%%%%%%%%%%%%%%%%%%%%%%%%%%%%%%%%%%%%%%%%%%%%%%%%%%%%%%%%%%%%%%%%%%%%%%%%%%%%%%
\subsection{The Classical Collapse Problem}
%%%%%%%%%%%%%%%%%%%%%%%%%%%%%%%%%%%%%%%%%%%%%%%%%%%%%%%%%%%%%%%%%%%%%%%%%%%%%%%%
Consider an electron of mass $m$ and charge $-e$ moving in the Coulomb
potential
\begin{equation}
V(r)=-\frac{e^2}{r},
\end{equation}
(Gaussian units).
Ignoring radiation reaction, Newton's equation is
\begin{equation}
m\ddot{\mathbf r}=-\frac{e^2}{r^3}\mathbf r.
\label{eq:CoulombNewton}
\end{equation}
Its solutions are the familiar Kepler ellipses,
\[E=-\frac{e^2}{2a},\]
where $a$ denotes the semi-major axis.
Within purely mechanical classical physics these orbits are perfectly
stable.
%%%%%%%%%%%%%%%%%%%%%%%%%%%%%%%%%%%%%%%%%%%%%%%%%%%%%%%%%%%%%%%%%%%%%%%%%%%%%%%%
\subsection{Instability Due to Radiation}
%%%%%%%%%%%%%%%%%%%%%%%%%%%%%%%%%%%%%%%%%%%%%%%%%%%%%%%%%%%%%%%%%%%%%%%%%%%%%%%%
The situation changes once Maxwell's equations are included.
Since the electron accelerates continuously, it radiates energy
according to the Larmor formula,
\begin{equation}
P=\frac{2e^2}{3c^3}a^2.
\label{eq:Larmor}
\end{equation}
The orbital energy therefore decreases,
\[\frac{dE}{dt}=-P,\]
causing the orbit to shrink.
A simple estimate predicts that a classical hydrogen atom would collapse
into the nucleus in approximately
\[10^{-11}\ {\rm s},\]
far shorter than any observed atomic lifetime.
This classical catastrophe was one of the principal motivations for the
development of quantum mechanics.
%%%%%%%%%%%%%%%%%%%%%%%%%%%%%%%%%%%%%%%%%%%%%%%%%%%%%%%%%%%%%%%%%%%%%%%%%%%%%%%%
\subsection{SED Perspective}
%%%%%%%%%%%%%%%%%%%%%%%%%%%%%%%%%%%%%%%%%%%%%%%%%%%%%%%%%%%%%%%%%%%%%%%%%%%%%%%%
Stochastic electrodynamics modifies this picture in a fundamental way.
The electron no longer evolves in the Coulomb field alone.
Instead,
\begin{equation}
m\ddot{\mathbf r}=-\frac{e^2}{r^3}\mathbf r+e\mathbf E_{\rm ZP}(t)+
    \mathbf F_{\rm RR},
\label{eq:HydrogenSED}
\end{equation}
where
\[\mathbf E_{\rm ZP}\]
is the stochastic zero-point field and
\[\mathbf F_{\rm RR}\]
denotes radiation reaction.
The dynamics now contains two competing mechanisms.
Radiation reaction continuously removes orbital energy,
while stochastic vacuum fluctuations continuously replenish it.
The fundamental issue is whether these two processes admit a stationary
statistical equilibrium.
%%%%%%%%%%%%%%%%%%%%%%%%%%%%%%%%%%%%%%%%%%%%%%%%%%%%%%%%%%%%%%%%%%%%%%%%%%%%%%%%
\subsection{Dimensionless Parameters}
%%%%%%%%%%%%%%%%%%%%%%%%%%%%%%%%%%%%%%%%%%%%%%%%%%%%%%%%%%%%%%%%%%%%%%%%%%%%%%%%
Several characteristic scales govern the problem.
The Bohr radius,
\begin{equation}
a_0=\frac{\hbar^2}{me^2},
\end{equation}
sets the natural length scale.
The corresponding energy is
\begin{equation}
E_{\rm B}=-\frac{me^4}{2\hbar^2},
\end{equation}
while the orbital frequency is approximately
\begin{equation}
\omega_{\rm B}=\frac{me^4}{\hbar^3}.
\end{equation}
Although these quantities contain Planck's constant, SED interprets
$\hbar$ merely as the normalization of the zero-point spectrum.
Consequently, the appearance of atomic scales is viewed as emerging from
the interaction between classical Coulomb dynamics and the stochastic
vacuum.
%%%%%%%%%%%%%%%%%%%%%%%%%%%%%%%%%%%%%%%%%%%%%%%%%%%%%%%%%%%%%%%%%%%%%%%%%%%%%%%%
\subsection{Mathematical Difficulties}
%%%%%%%%%%%%%%%%%%%%%%%%%%%%%%%%%%%%%%%%%%%%%%%%%%%%%%%%%%%%%%%%%%%%%%%%%%%%%%%%
Compared with the harmonic oscillator, the hydrogen atom introduces
several new mathematical complications.
\begin{enumerate}
\item
The Coulomb force is nonlinear.
\item
The orbital frequency is not constant but depends on the instantaneous
trajectory.
\item
Highly eccentric orbits contain many Fourier harmonics.
\item
Radiation reaction becomes strongly trajectory dependent.
\item
The stochastic force excites infinitely many resonances.
\item
Close encounters with the nucleus generate very large accelerations.
\end{enumerate}
These difficulties prevent an exact analytical solution.
Consequently, nearly all investigations rely either on perturbation
theory or on extensive numerical simulations.
%%%%%%%%%%%%%%%%%%%%%%%%%%%%%%%%%%%%%%%%%%%%%%%%%%%%%%%%%%%%%%%%%%%%%%%%%%%%%%%%
\subsection{The Central Physical Question}
%%%%%%%%%%%%%%%%%%%%%%%%%%%%%%%%%%%%%%%%%%%%%%%%%%%%%%%%%%%%%%%%%%%%%%%%%%%%%%%%
The harmonic oscillator possesses a unique stable equilibrium.
The Coulomb problem is fundamentally different.
Instead of asking whether the particle fluctuates about a fixed point,
one must determine whether the probability distribution in phase space
approaches a stationary measure.
In modern language the problem becomes one of stochastic dynamics on the
Kepler phase space.
Does there exist an invariant probability measure
\[\rho(\mathbf r,\mathbf p)\]
satisfying
\[\frac{\partial\rho}{\partial t}=0?\]
If such a measure exists, one may compare it directly with the quantum
ground-state distribution of hydrogen.
%%%%%%%%%%%%%%%%%%%%%%%%%%%%%%%%%%%%%%%%%%%%%%%%%%%%%%%%%%%%%%%%%%%%%%%%%%%%%%%%
\subsection*{Editorial Comment}
%%%%%%%%%%%%%%%%%%%%%%%%%%%%%%%%%%%%%%%%%%%%%%%%%%%%%%%%%%%%%%%%%%%%%%%%%%%%%%%%
Historically, discussions of SED often begin immediately with numerical
simulations or with particular analytical models. We believe that this
obscures the essential physical issue.
The hydrogen atom is fundamentally a problem of dynamical balance
between two competing mechanisms: deterministic radiative energy loss
and stochastic energy injection from the electromagnetic vacuum. The
existence, uniqueness, and stability of the resulting statistical
equilibrium constitute the central mathematical questions.
From this perspective the hydrogen atom is not merely another example of
SED but rather the decisive test of the entire research programme.
%%%%%%%%%%%%%%%%%%%%%%%%%%%%%%%%%%%%%%%%%%%%%%%%%%%%%%%%%%%%%%%%%%%%%%%%%%%%%%%%
\subsection{Analytical Approaches: Puthoff's Energy-Balance Model}
\label{sec:Puthoff}
%%%%%%%%%%%%%%%%%%%%%%%%%%%%%%%%%%%%%%%%%%%%%%%%%%%%%%%%%%%%%%%%%%%%%%%%%%%%%%%%
The first widely discussed analytical explanation of atomic stability
within stochastic electrodynamics was proposed by Theodore Puthoff in
1987. Rather than attempting to solve the complete stochastic equations
of motion, Puthoff considered the average balance between radiative
energy loss and energy absorbed from the classical zero-point field.
The basic hypothesis is simple.
A stationary atomic orbit is possible if the average power radiated by
the electron equals the average power absorbed from the stochastic
vacuum,
\begin{equation}
\boxed{\left\langle P_{\rm rad}\right\rangle=
\left\langle P_{\rm abs}\right\rangle .}
\label{eq:EnergyBalance}
\end{equation}
This condition is analogous to thermal equilibrium in statistical
physics and represents the central physical idea of Puthoff's model.
%%%%%%%%%%%%%%%%%%%%%%%%%%%%%%%%%%%%%%%%%%%%%%%%%%%%%%%%%%%%%%%%%%%%%%%%%%%%%%%%
\subsubsection{Radiated Power}
%%%%%%%%%%%%%%%%%%%%%%%%%%%%%%%%%%%%%%%%%%%%%%%%%%%%%%%%%%%%%%%%%%%%%%%%%%%%%%%%
For a nonrelativistic charge moving with acceleration
$\mathbf a$, the Larmor formula gives
\begin{equation}
P_{\rm rad}=\frac{2e^2}{3c^3}a^2.
\label{eq:LarmorAgain}
\end{equation}
For a circular Coulomb orbit,
\begin{equation}
\frac{mv^2}{r}=\frac{e^2}{r^2},
\end{equation}
so that
\begin{equation}
a=\frac{v^2}{r}=\frac{e^2}{mr^2}.
\end{equation}
Substitution into Eq.~(\ref{eq:LarmorAgain}) yields
\begin{equation}
P_{\rm rad}=\frac{2e^6}{3m^2c^3r^4}.
\label{eq:PRadCircular}
\end{equation}
Thus radiative losses increase rapidly as the orbital radius decreases.
%%%%%%%%%%%%%%%%%%%%%%%%%%%%%%%%%%%%%%%%%%%%%%%%%%%%%%%%%%%%%%%%%%%%%%%%%%%%%%%%
\subsubsection{Energy Absorption}
%%%%%%%%%%%%%%%%%%%%%%%%%%%%%%%%%%%%%%%%%%%%%%%%%%%%%%%%%%%%%%%%%%%%%%%%%%%%%%%%
The interaction with the stochastic vacuum continuously transfers energy
to the orbiting electron.
Puthoff estimated the average absorbed power by treating the electron as
a resonant system coupled to the zero-point field. The calculation
employs the Lorentz-invariant spectral density
\begin{equation}
u(\omega)=\frac{\hbar\omega^3}{2\pi^2c^3},
\end{equation}
together with the dipole coupling between the electron and the
electromagnetic field.
The resulting absorbed power may be written schematically as
\begin{equation}
P_{\rm abs}=\frac{e^2\hbar\omega^3}{2mc^3},
\label{eq:PAbs}
\end{equation}
where $\omega$ is the orbital frequency. A more careful derivation,
including angular averaging and polarization sums, leads to the same
functional dependence.
%%%%%%%%%%%%%%%%%%%%%%%%%%%%%%%%%%%%%%%%%%%%%%%%%%%%%%%%%%%%%%%%%%%%%%%%%%%%%%%%
\subsubsection{Orbital Frequency}
%%%%%%%%%%%%%%%%%%%%%%%%%%%%%%%%%%%%%%%%%%%%%%%%%%%%%%%%%%%%%%%%%%%%%%%%%%%%%%%%
For a circular Kepler orbit,
\begin{equation}
m\omega^2r=\frac{e^2}{r^2},
\end{equation}
which implies
\begin{equation}
\omega=\sqrt{\frac{e^2}{mr^3}}.
\label{eq:KeplerFrequency}
\end{equation}
The absorbed power therefore scales as
\begin{equation}
P_{\rm abs}\propto\frac{e^5\hbar}{m^{5/2}c^3r^{9/2}}.
\end{equation}
%%%%%%%%%%%%%%%%%%%%%%%%%%%%%%%%%%%%%%%%%%%%%%%%%%%%%%%%%%%%%%%%%%%%%%%%%%%%%%%%
\subsubsection{Stationary Radius}
%%%%%%%%%%%%%%%%%%%%%%%%%%%%%%%%%%%%%%%%%%%%%%%%%%%%%%%%%%%%%%%%%%%%%%%%%%%%%%%%
The equilibrium condition,
\[P_{\rm rad}=P_{\rm abs},\]
determines a characteristic orbital radius.
Carrying out the algebra gives
\begin{equation}
r\sim\frac{\hbar^2}{me^2},
\end{equation}
which is precisely the Bohr radius,
\begin{equation}
\boxed{a_0=\frac{\hbar^2}{me^2}.}
\end{equation}
Within SED this result is interpreted not as a consequence of angular
momentum quantization but as the radius at which energy loss through
radiation is exactly compensated by energy gained from the stochastic
vacuum.
%%%%%%%%%%%%%%%%%%%%%%%%%%%%%%%%%%%%%%%%%%%%%%%%%%%%%%%%%%%%%%%%%%%%%%%%%%%%%%%%
\subsubsection{Interpretation}
%%%%%%%%%%%%%%%%%%%%%%%%%%%%%%%%%%%%%%%%%%%%%%%%%%%%%%%%%%%%%%%%%%%%%%%%%%%%%%%%
The significance of Puthoff's calculation lies less in the numerical
value of the Bohr radius than in the underlying physical picture.
In the Bohr model, the stationary orbit is postulated.
In Schrödinger's theory, the ground state is obtained as an eigenstate
of the Hamiltonian.
In stochastic electrodynamics, the characteristic atomic scale is
interpreted as a nonequilibrium steady state maintained by the continual
exchange of energy between the electron and the vacuum field.
The atom is therefore viewed as a dynamical system rather than a static
bound state.
%%%%%%%%%%%%%%%%%%%%%%%%%%%%%%%%%%%%%%%%%%%%%%%%%%%%%%%%%%%%%%%%%%%%%%%%%%%%%%%%
\subsubsection*{Critical Assessment}
%%%%%%%%%%%%%%%%%%%%%%%%%%%%%%%%%%%%%%%%%%%%%%%%%%%%%%%%%%%%%%%%%%%%%%%%%%%%%%%%
Puthoff's model represented an important milestone because it showed
that the observed atomic length scale can emerge from a balance between
classical radiation reaction and stochastic vacuum fluctuations.
Nevertheless, the derivation relies on a number of simplifying
assumptions.
First, the analysis assumes circular orbits, whereas the Coulomb problem
admits a continuous family of elliptical trajectories with varying
eccentricities. Second, the absorbed power is obtained using approximate
response functions rather than from an exact solution of the stochastic
equations of motion. Third, the equilibrium condition is imposed on
ensemble-averaged powers and does not by itself establish the existence
or stability of a stationary probability distribution in phase space.
Consequently, Puthoff's calculation should be regarded as an important
consistency argument rather than as a complete derivation of the
hydrogen ground state.
From a modern perspective, the principal contribution of this work is
conceptual: it demonstrated that the characteristic atomic scales of
hydrogen are not obviously incompatible with a classical stochastic
vacuum and thereby motivated the more sophisticated numerical
investigations discussed in the following sections.
%%%%%%%%%%%%%%%%%%%%%%%%%%%%%%%%%%%%%%%%%%%%%%%%%%%%%%%%%%%%%%%%%%%%%%%%%%%%%%%%
\subsection{A Modern Dynamical Systems Formulation}
\label{sec:SEDModern}
%%%%%%%%%%%%%%%%%%%%%%%%%%%%%%%%%%%%%%%%%%%%%%%%%%%%%%%%%%%%%%%%%%%%%%%%%%%%%%%%
The analytical models discussed in the previous section provide valuable
physical insight but leave unanswered an important mathematical
question. An equilibrium condition based on the average balance between
absorbed and radiated power does not, by itself, establish the existence
of a statistically stationary atomic state.
From the viewpoint of modern stochastic dynamics, the central question
is different.
Rather than asking whether the average energy remains constant, one asks
whether the stochastic evolution possesses an invariant probability
measure on phase space.
%%%%%%%%%%%%%%%%%%%%%%%%%%%%%%%%%%%%%%%%%%%%%%%%%%%%%%%%%%%%%%%%%%%%%%%%%%%%%%%%
\subsubsection{Phase-Space Dynamics}
%%%%%%%%%%%%%%%%%%%%%%%%%%%%%%%%%%%%%%%%%%%%%%%%%%%%%%%%%%%%%%%%%%%%%%%%%%%%%%%%
Let
\[X(t)=({\bf r}(t),{\bf p}(t))\]
denote the six-dimensional phase-space point of the electron.
The deterministic Coulomb dynamics is generated by the Hamiltonian
\begin{equation}
H({\bf r},{\bf p})=\frac{{\bf p}^2}{2m}-\frac{e^2}{r}.
\end{equation}
In stochastic electrodynamics the equations of motion become
\begin{align}
\dot{\bf r}&=\frac{\partial H}{\partial{\bf p}},\\
\dot{\bf p}&=-\frac{\partial H}{\partial{\bf r}}+
{\bf F}_{\rm RR}+e{\bf E}_{\rm ZP}(t),
\label{eq:SEDPhaseSpace}
\end{align}
where
\[{\bf F}_{\rm RR}\]
is the radiation-reaction force and
\[{\bf E}_{\rm ZP}(t)\]
is the stochastic zero-point field.
The deterministic Hamiltonian flow is therefore replaced by a random
flow on phase space.
%%%%%%%%%%%%%%%%%%%%%%%%%%%%%%%%%%%%%%%%%%%%%%%%%%%%%%%%%%%%%%%%%%%%%%%%%%%%%%%%
\subsubsection{Non-Markovian Character}
%%%%%%%%%%%%%%%%%%%%%%%%%%%%%%%%%%%%%%%%%%%%%%%%%%%%%%%%%%%%%%%%%%%%%%%%%%%%%%%%
Unlike Nelson's stochastic mechanics, the random force is not assumed to
be white noise.
Instead,
\[\langle E_i(t)E_j(t')\rangle=C_{ij}(t-t'),\]
where the correlation function is determined by the zero-point spectrum.
The stochastic forcing therefore possesses memory.
Consequently,
the electron dynamics is generally non-Markovian.
This observation has important mathematical consequences.
Many standard results concerning diffusion processes, Fokker--Planck
equations and invariant measures cannot be applied directly.
%%%%%%%%%%%%%%%%%%%%%%%%%%%%%%%%%%%%%%%%%%%%%%%%%%%%%%%%%%%%%%%%%%%%%%%%%%%%%%%%
\subsubsection{Extended Markov Description}
%%%%%%%%%%%%%%%%%%%%%%%%%%%%%%%%%%%%%%%%%%%%%%%%%%%%%%%%%%%%%%%%%%%%%%%%%%%%%%%%
A common strategy in stochastic dynamics is to enlarge the state space.
Instead of considering only
\[({\bf r},{\bf p}),\]
one introduces additional variables describing the stochastic field.
The enlarged process
\[({\bf r},{\bf p},{\bf E}_{\rm ZP})\]
is then Markovian even though the reduced particle dynamics is not.
This construction is widely used in the theory of colored-noise
processes and generalized Langevin equations.
It provides a natural mathematical framework for stochastic
electrodynamics.
%%%%%%%%%%%%%%%%%%%%%%%%%%%%%%%%%%%%%%%%%%%%%%%%%%%%%%%%%%%%%%%%%%%%%%%%%%%%%%%%
\subsubsection{Invariant Probability Measure}
%%%%%%%%%%%%%%%%%%%%%%%%%%%%%%%%%%%%%%%%%%%%%%%%%%%%%%%%%%%%%%%%%%%%%%%%%%%%%%%%
The principal mathematical problem may now be formulated precisely.
Does there exist a probability density
\[\rho({\bf r},{\bf p})\]
such that
\begin{equation}
\frac{\partial\rho}{\partial t}=0
\end{equation}
under the stochastic dynamics generated by
Eq.~(\ref{eq:SEDPhaseSpace})?
If such a measure exists, several further questions arise.
\begin{itemize}
\item
Is it unique?
\item
Is it normalizable?
\item
Is it stable against perturbations?
\item
Is it ergodic?
\item
Does it coincide with the quantum ground-state distribution?
\end{itemize}
These questions are considerably stronger than the simple energy-balance
condition discussed previously.
%%%%%%%%%%%%%%%%%%%%%%%%%%%%%%%%%%%%%%%%%%%%%%%%%%%%%%%%%%%%%%%%%%%%%%%%%%%%%%%%
\subsubsection{Generator of the Stochastic Dynamics}
%%%%%%%%%%%%%%%%%%%%%%%%%%%%%%%%%%%%%%%%%%%%%%%%%%%%%%%%%%%%%%%%%%%%%%%%%%%%%%%%
If the stochastic process were Markovian, its evolution would be
generated by an operator
\[L,\]
acting on smooth functions of phase space,
\[\frac{d}{dt}\langle f\rangle=\langle Lf\rangle.\]
The invariant distribution would satisfy
\begin{equation}
L^*\rho=0,
\end{equation}
where
\[L^*\]
is the adjoint generator.
For white-noise processes this reduces to the familiar
Fokker--Planck equation.
For SED, however, the colored nature of the zero-point field generally
leads to generalized kinetic equations containing memory kernels.
%%%%%%%%%%%%%%%%%%%%%%%%%%%%%%%%%%%%%%%%%%%%%%%%%%%%%%%%%%%%%%%%%%%%%%%%%%%%%%%%
\subsubsection{Relation to Energy Balance}
%%%%%%%%%%%%%%%%%%%%%%%%%%%%%%%%%%%%%%%%%%%%%%%%%%%%%%%%%%%%%%%%%%%%%%%%%%%%%%%%
The energy-balance condition may now be interpreted as a necessary but
not sufficient condition for stationarity.
Indeed,
if a stationary probability distribution exists,
\[\frac{d}{dt}\langle H\rangle=0.\]
Therefore,
\[\left<P_{\rm abs}\right>=\left<P_{\rm rad}\right>.\]
The converse implication does not necessarily hold.
Equality of average powers does not guarantee the existence,
uniqueness, or stability of an invariant probability measure.
%%%%%%%%%%%%%%%%%%%%%%%%%%%%%%%%%%%%%%%%%%%%%%%%%%%%%%%%%%%%%%%%%%%%%%%%%%%%%%%%
\subsubsection*{Critical Assessment}
%%%%%%%%%%%%%%%%%%%%%%%%%%%%%%%%%%%%%%%%%%%%%%%%%%%%%%%%%%%%%%%%%%%%%%%%%%%%%%%%
Reformulating stochastic electrodynamics in terms of stochastic
dynamical systems clarifies several conceptual issues. It distinguishes
between local energy balance and global statistical equilibrium and
provides a natural framework for comparing SED with other stochastic
approaches discussed in this monograph.
At the same time, the reformulation exposes a major mathematical
challenge. Because the zero-point field possesses long-range temporal
correlations, the particle dynamics is generally non-Markovian.
Rigorous existence and uniqueness theorems for invariant measures are
therefore considerably more difficult to establish than for ordinary
diffusion processes.
This observation suggests that future developments of SED may benefit
from modern techniques in the theory of generalized Langevin equations,
hypoelliptic diffusions, random dynamical systems, and stochastic
Hamiltonian mechanics. Such methods have matured substantially since the
original formulation of SED and may provide new insight into the
question of atomic stability.
%%%%%%%%%%%%%%%%%%%%%%%%%%%%%%%%%%%%%%%%%%%%%%%%%%%%%%%%%%%%%%%%%%%%%%%%%%%%%%%%
\subsection{Numerical Simulations: The Work of Cole and Zou}
\label{sec:ColeZou}
%%%%%%%%%%%%%%%%%%%%%%%%%%%%%%%%%%%%%%%%%%%%%%%%%%%%%%%%%%%%%%%%%%%%%%%%%%%%%%%%
The analytical models discussed in the previous sections rely on various
approximations concerning the interaction between the electron and the
stochastic vacuum field. Because of the complexity of the full
stochastic equations, an important alternative is direct numerical
integration.
The first widely discussed numerical investigation was carried out by
Cole and Zou in 2003. Their objective was straightforward: integrate the
classical equations of motion of an electron in the Coulomb field,
including both radiation reaction and the stochastic zero-point
electromagnetic field, and determine whether the long-time statistical
behavior resembles the quantum-mechanical ground state.
%%%%%%%%%%%%%%%%%%%%%%%%%%%%%%%%%%%%%%%%%%%%%%%%%%%%%%%%%%%%%%%%%%%%%%%%%%%%%%%%
\subsubsection{Equations of Motion}
%%%%%%%%%%%%%%%%%%%%%%%%%%%%%%%%%%%%%%%%%%%%%%%%%%%%%%%%%%%%%%%%%%%%%%%%%%%%%%%%
The electron dynamics was modeled by
\begin{equation}
m\ddot{\mathbf r}=-\frac{e^2}{r^3}{\mathbf r}+e{\mathbf E}_{\rm ZP}(t)
    +{\mathbf F}_{\rm RR},
\label{eq:ColeZouEq}
\end{equation}
where
\begin{itemize}
\item the first term is the Coulomb force,
\item the second term represents the stochastic zero-point field,
\item the third term accounts for radiation reaction.
\end{itemize}
The stochastic field was represented by a finite but very large sum of
plane-wave modes,
\begin{equation}
{\mathbf E}_{\rm ZP}(t)=\sum_{n=1}^{N}{\mathbf E}_n\cos(\omega_nt+\phi_n),
\end{equation}
where the random phases
\[\phi_n\]
were chosen independently and uniformly on
\[[0,2\pi),\]
while the amplitudes were adjusted to reproduce the Lorentz-invariant
zero-point spectrum discussed earlier.
%%%%%%%%%%%%%%%%%%%%%%%%%%%%%%%%%%%%%%%%%%%%%%%%%%%%%%%%%%%%%%%%%%%%%%%%%%%%%%%%
\subsubsection{Finite Mode Representation}
%%%%%%%%%%%%%%%%%%%%%%%%%%%%%%%%%%%%%%%%%%%%%%%%%%%%%%%%%%%%%%%%%%%%%%%%%%%%%%%%
An exact realization of the zero-point field requires infinitely many
electromagnetic modes. Numerical simulations necessarily replace this by
a finite expansion,
\[N<\infty.\]
This approximation introduces two artificial scales:
\begin{itemize}
\item a lowest represented frequency,
\item a highest represented frequency.
\end{itemize}
Consequently, every simulation implicitly introduces infrared and
ultraviolet cutoffs.
The dependence of the numerical results upon these cutoffs is therefore
an important issue.
%%%%%%%%%%%%%%%%%%%%%%%%%%%%%%%%%%%%%%%%%%%%%%%%%%%%%%%%%%%%%%%%%%%%%%%%%%%%%%%%
\subsubsection{Integration Algorithm}
%%%%%%%%%%%%%%%%%%%%%%%%%%%%%%%%%%%%%%%%%%%%%%%%%%%%%%%%%%%%%%%%%%%%%%%%%%%%%%%%
The equations of motion were integrated over many orbital periods using
high-order adaptive numerical methods.
Because the stochastic force contains frequencies extending far above
the orbital frequency, extremely small time steps are required in order
to avoid numerical instability.
Long integration times are essential since the objective is not a single
trajectory but the stationary statistical distribution obtained after
transients have decayed.
%%%%%%%%%%%%%%%%%%%%%%%%%%%%%%%%%%%%%%%%%%%%%%%%%%%%%%%%%%%%%%%%%%%%%%%%%%%%%%%%
\subsubsection{Observed Probability Distribution}
%%%%%%%%%%%%%%%%%%%%%%%%%%%%%%%%%%%%%%%%%%%%%%%%%%%%%%%%%%%%%%%%%%%%%%%%%%%%%%%%
The principal observable is the radial probability density,
\[P(r).\]
Cole and Zou reported that the numerically obtained distribution
approaches a form close to the quantum-mechanical ground-state density,
\begin{equation}
P_{\rm QM}(r)=\frac{4r^2}{a_0^3}e^{-2r/a_0},
\end{equation}
where
\[a_0=\frac{\hbar^2}{me^2}\]
is the Bohr radius.
Although the agreement was not exact, the qualitative similarity was
remarkable given that the underlying dynamics remained entirely
classical.
%%%%%%%%%%%%%%%%%%%%%%%%%%%%%%%%%%%%%%%%%%%%%%%%%%%%%%%%%%%%%%%%%%%%%%%%%%%%%%%%
\subsubsection{Energy Distribution}
%%%%%%%%%%%%%%%%%%%%%%%%%%%%%%%%%%%%%%%%%%%%%%%%%%%%%%%%%%%%%%%%%%%%%%%%%%%%%%%%
The simulations also indicated that the orbital energy fluctuates
continuously around the ground-state value rather than remaining fixed.
This behavior differs from the stationary quantum eigenstate, whose
energy is sharply defined in the absence of measurement.
Within SED, however, continuous stochastic exchange of energy with the
vacuum field is expected.
Consequently, fluctuations of the instantaneous orbital energy are a
natural feature of the theory.
%%%%%%%%%%%%%%%%%%%%%%%%%%%%%%%%%%%%%%%%%%%%%%%%%%%%%%%%%%%%%%%%%%%%%%%%%%%%%%%%
\subsubsection{Interpretation}
%%%%%%%%%%%%%%%%%%%%%%%%%%%%%%%%%%%%%%%%%%%%%%%%%%%%%%%%%%%%%%%%%%%%%%%%%%%%%%%%
The numerical results suggest the existence of a long-lived statistical
equilibrium maintained by the balance between stochastic excitation and
radiative damping.
The atom is therefore interpreted as a nonequilibrium steady state.
Individual trajectories remain highly irregular, but the ensemble of
trajectories approaches an approximately stationary probability
distribution.
This interpretation differs fundamentally from the quantum-mechanical
picture, where the ground state is an energy eigenstate described by a
stationary wave function.
%%%%%%%%%%%%%%%%%%%%%%%%%%%%%%%%%%%%%%%%%%%%%%%%%%%%%%%%%%%%%%%%%%%%%%%%%%%%%%%%
\subsubsection{Numerical Limitations}
%%%%%%%%%%%%%%%%%%%%%%%%%%%%%%%%%%%%%%%%%%%%%%%%%%%%%%%%%%%%%%%%%%%%%%%%%%%%%%%%
Several numerical approximations should be kept in mind.
\begin{enumerate}
\item
The stochastic vacuum was represented by finitely many modes.
\item
Finite integration time may prevent complete equilibration.
\item
Radiation reaction was treated approximately.
\item
Very high frequencies could not be included explicitly.
\item
The calculations were restricted primarily to the ground state.
\end{enumerate}
These limitations are common to virtually all large-scale SED
simulations and should not be interpreted as shortcomings specific to
the work of Cole and Zou.
%%%%%%%%%%%%%%%%%%%%%%%%%%%%%%%%%%%%%%%%%%%%%%%%%%%%%%%%%%%%%%%%%%%%%%%%%%%%%%%%
\subsubsection*{Critical Assessment}
%%%%%%%%%%%%%%%%%%%%%%%%%%%%%%%%%%%%%%%%%%%%%%%%%%%%%%%%%%%%%%%%%%%%%%%%%%%%%%%%
The simulations of Cole and Zou constitute one of the most significant
computational achievements in the history of stochastic
electrodynamics. They demonstrated that direct numerical integration of
the classical stochastic equations can produce probability
distributions resembling those of the quantum-mechanical hydrogen ground
state.
At the same time, the conclusions must be interpreted with appropriate
caution. The simulations employed finite spectral representations of
the vacuum field, finite integration times, and approximate treatments
of radiation reaction. Moreover, agreement with the ground-state radial
distribution does not automatically imply agreement with all quantum
observables.
Nevertheless, the work changed the direction of research. Before 2003,
the discussion of atomic stability within SED was dominated primarily by
analytical arguments. After Cole and Zou, the emphasis shifted toward
large-scale numerical investigations, opening the possibility of
testing the theory quantitatively rather than relying solely on
heuristic reasoning.
%%%%%%%%%%%%%%%%%%%%%%%%%%%%%%%%%%%%%%%%%%%%%%%%%%%%%%%%%%%%%%%%%%%%%%%%%%%%%%%%
\subsection{Large-Scale Numerical Investigations:
Nieuwenhuizen and Liska}
\label{sec:Nieuwenhuizen}
%%%%%%%%%%%%%%%%%%%%%%%%%%%%%%%%%%%%%%%%%%%%%%%%%%%%%%%%%%%%%%%%%%%%%%%%%%%%%%%%
Following the pioneering work of Cole and Zou, substantially larger and
more accurate numerical investigations were undertaken by Theo M.
Nieuwenhuizen, Richard Liska and collaborators. Their objective was not
merely to reproduce earlier simulations but to determine whether the
hydrogen atom remains statistically stable over integration times and
spectral resolutions far exceeding those previously considered.
These investigations constitute the most extensive numerical study of
stochastic electrodynamics performed to date and have significantly
clarified both the strengths and the limitations of the theory.
%%%%%%%%%%%%%%%%%%%%%%%%%%%%%%%%%%%%%%%%%%%%%%%%%%%%%%%%%%%%%%%%%%%%%%%%%%%%%%%%
\subsubsection{Motivation}
%%%%%%%%%%%%%%%%%%%%%%%%%%%%%%%%%%%%%%%%%%%%%%%%%%%%%%%%%%%%%%%%%%%%%%%%%%%%%%%%
The calculations of Cole and Zou demonstrated encouraging agreement
with the quantum-mechanical ground-state distribution. However, several
questions remained unanswered.
\begin{itemize}
\item
Does the apparent equilibrium persist over much longer times?
\item
Is the stationary distribution independent of the numerical
discretization?
\item
How sensitive are the results to the ultraviolet cutoff?
\item
Does the stochastic dynamics eventually produce self-ionization?
\item
Are the observed probability distributions genuinely stationary?
\end{itemize}
Answering these questions requires simulations extending over many more
orbital periods and involving significantly larger numbers of
electromagnetic modes.
%%%%%%%%%%%%%%%%%%%%%%%%%%%%%%%%%%%%%%%%%%%%%%%%%%%%%%%%%%%%%%%%%%%%%%%%%%%%%%%%
\subsubsection{Numerical Improvements}
%%%%%%%%%%%%%%%%%%%%%%%%%%%%%%%%%%%%%%%%%%%%%%%%%%%%%%%%%%%%%%%%%%%%%%%%%%%%%%%%
Compared with earlier work, Nieuwenhuizen and Liska introduced several
important improvements.
These include
\begin{itemize}
\item
larger spectral representations of the zero-point field,
\item
improved treatment of orbital eccentricity,
\item
adaptive numerical integration,
\item
longer simulation times,
\item
careful monitoring of conserved quantities,
\item
statistical averaging over many independent realizations.
\end{itemize}
The resulting computations required computational resources
substantially exceeding those available for the original Cole--Zou
study.
%%%%%%%%%%%%%%%%%%%%%%%%%%%%%%%%%%%%%%%%%%%%%%%%%%%%%%%%%%%%%%%%%%%%%%%%%%%%%%%%
\subsubsection{Observed Dynamics}
%%%%%%%%%%%%%%%%%%%%%%%%%%%%%%%%%%%%%%%%%%%%%%%%%%%%%%%%%%%%%%%%%%%%%%%%%%%%%%%%
During the early stages of the simulations the electron generally
approaches a statistical state whose radial distribution resembles the
quantum-mechanical ground state.
However, over sufficiently long integration times a qualitatively
different behavior may emerge.
Rare but large stochastic fluctuations can transfer sufficient energy
to the electron to drive it into highly eccentric trajectories. Under
certain conditions these excursions are followed by complete
self-ionization, after which the electron escapes permanently from the
Coulomb potential.
The occurrence of such events depends sensitively on the numerical
implementation and the representation of the stochastic vacuum.
%%%%%%%%%%%%%%%%%%%%%%%%%%%%%%%%%%%%%%%%%%%%%%%%%%%%%%%%%%%%%%%%%%%%%%%%%%%%%%%%
\subsubsection{Metastability Versus Stability}
%%%%%%%%%%%%%%%%%%%%%%%%%%%%%%%%%%%%%%%%%%%%%%%%%%%%%%%%%%%%%%%%%%%%%%%%%%%%%%%%
The simulations suggest an important distinction between two concepts
that are often conflated.
A system may possess a long-lived metastable distribution without
possessing a strictly stationary invariant measure.
From the perspective of stochastic dynamics, the relevant question is
therefore whether the hydrogen atom in SED is
\begin{enumerate}
\item
strictly stationary,
\item
metastable with an extremely long lifetime,
\item
or ultimately unstable.
\end{enumerate}
The available numerical evidence has not yet produced a universally
accepted answer.
%%%%%%%%%%%%%%%%%%%%%%%%%%%%%%%%%%%%%%%%%%%%%%%%%%%%%%%%%%%%%%%%%%%%%%%%%%%%%%%%
\subsubsection{Role of High-Frequency Modes}
%%%%%%%%%%%%%%%%%%%%%%%%%%%%%%%%%%%%%%%%%%%%%%%%%%%%%%%%%%%%%%%%%%%%%%%%%%%%%%%%
One of the principal lessons of the large-scale simulations concerns the
importance of the ultraviolet portion of the zero-point spectrum.
Because the spectral density grows as
\[u(\omega)\propto\omega^3,\]
high-frequency modes carry an increasingly large amount of energy.
Although their direct coupling to the orbital motion is weak, nonlinear
interactions and harmonic generation can transfer energy across widely
separated frequency scales.
Consequently, numerical truncation of the spectrum may influence the
long-time dynamics in subtle ways.
%%%%%%%%%%%%%%%%%%%%%%%%%%%%%%%%%%%%%%%%%%%%%%%%%%%%%%%%%%%%%%%%%%%%%%%%%%%%%%%%
\subsubsection{Numerical Convergence}
%%%%%%%%%%%%%%%%%%%%%%%%%%%%%%%%%%%%%%%%%%%%%%%%%%%%%%%%%%%%%%%%%%%%%%%%%%%%%%%%
A recurring difficulty concerns the notion of convergence.
Increasing the number of electromagnetic modes changes not only the
accuracy of the stochastic field but also the physical fluctuations
experienced by the electron.
Unlike ordinary discretization errors, spectral refinement therefore
modifies the effective stochastic dynamics itself.
Demonstrating convergence requires showing that observable probability
distributions become independent of
\begin{itemize}
\item
spectral resolution,
\item
time step,
\item
integration algorithm,
\item
simulation duration,
\item
initial conditions.
\end{itemize}
Establishing all of these properties simultaneously remains a demanding
computational problem.
%%%%%%%%%%%%%%%%%%%%%%%%%%%%%%%%%%%%%%%%%%%%%%%%%%%%%%%%%%%%%%%%%%%%%%%%%%%%%%%%
\subsubsection{Interpretation}
%%%%%%%%%%%%%%%%%%%%%%%%%%%%%%%%%%%%%%%%%%%%%%%%%%%%%%%%%%%%%%%%%%%%%%%%%%%%%%%%
The work of Nieuwenhuizen and Liska shifted the discussion of SED away
from the simple question
\begin{quote}
``Does the hydrogen atom remain stable?''
\end{quote}
toward a more refined set of mathematical questions concerning
stochastic stability, invariant measures, escape probabilities and
rare-event statistics.
This change of viewpoint represents an important conceptual advance.
%%%%%%%%%%%%%%%%%%%%%%%%%%%%%%%%%%%%%%%%%%%%%%%%%%%%%%%%%%%%%%%%%%%%%%%%%%%%%%%%
\subsubsection*{Critical Assessment}
%%%%%%%%%%%%%%%%%%%%%%%%%%%%%%%%%%%%%%%%%%%%%%%%%%%%%%%%%%%%%%%%%%%%%%%%%%%%%%%%
The numerical investigations of Nieuwenhuizen and Liska constitute the
most thorough computational examination of the hydrogen atom within
stochastic electrodynamics currently available. Their work confirmed
that the stochastic dynamics can reproduce many qualitative features of
the quantum ground state over extended periods while simultaneously
revealing previously unrecognized challenges associated with long-time
evolution and ultraviolet fluctuations.
Perhaps the most significant contribution of these studies is not a
definitive answer regarding atomic stability but rather a clearer
formulation of the problem itself. The question is no longer whether
short numerical trajectories resemble quantum mechanics; instead, it is
whether the stochastic dynamics admits a robust invariant probability
measure under systematic refinement of the numerical model.
From the perspective of modern stochastic dynamics, this reformulation
is arguably the most important outcome of the large-scale simulations.
%%%%%%%%%%%%%%%%%%%%%%%%%%%%%%%%%%%%%%%%%%%%%%%%%%%%%%%%%%%%%%%%%%%%%%%%%%%%%%%%
\section{Numerical Reliability and Verification of SED Simulations}
\label{sec:NumericalReliability}
%%%%%%%%%%%%%%%%%%%%%%%%%%%%%%%%%%%%%%%%%%%%%%%%%%%%%%%%%%%%%%%%%%%%%%%%%%%%%%%%
The numerical investigation of stochastic electrodynamics presents
challenges that differ substantially from those encountered in ordinary
deterministic dynamical systems. The objective is not merely to compute
individual trajectories but to determine statistical properties of a
high-dimensional stochastic dynamical system over extremely long time
intervals.
Consequently, numerical accuracy alone is insufficient. One must also
establish statistical convergence, robustness with respect to numerical
parameters, and independence of arbitrary approximations introduced in
the representation of the stochastic vacuum.
This section summarizes the principal numerical issues that must be
addressed before conclusions concerning atomic stability can be regarded
as reliable.
%%%%%%%%%%%%%%%%%%%%%%%%%%%%%%%%%%%%%%%%%%%%%%%%%%%%%%%%%%%%%%%%%%%%%%%%%%%%%%%%
\subsection{Representation of the Stochastic Vacuum}
%%%%%%%%%%%%%%%%%%%%%%%%%%%%%%%%%%%%%%%%%%%%%%%%%%%%%%%%%%%%%%%%%%%%%%%%%%%%%%%%
The Lorentz-invariant zero-point field formally contains infinitely many
electromagnetic modes,
\begin{equation}
{\bf E}({\bf r},t)=\sum_{\lambda=1}^{2}\int d^3k\,
{\bf e}_\lambda({\bf k})A_\lambda({\bf k})
\cos({\bf k}\cdot{\bf r}-\omega t+\phi_\lambda).
\label{eq:InfiniteField}
\end{equation}
Direct numerical evaluation of this expression is impossible.
Every simulation therefore replaces the continuum by a finite expansion,
\begin{equation}
{\bf E}_N({\bf r},t)=\sum_{n=1}^{N}{\bf E}_n\cos(\omega_nt+\phi_n).
\label{eq:FiniteExpansion}
\end{equation}
The first numerical question is therefore obvious:
\begin{quote}
How large must $N$ be before observable quantities become independent
of further increases?
\end{quote}
Unlike ordinary Fourier approximations, increasing $N$ changes the
physical stochastic forcing itself.
Consequently, convergence cannot be inferred solely from pointwise field
errors.
%%%%%%%%%%%%%%%%%%%%%%%%%%%%%%%%%%%%%%%%%%%%%%%%%%%%%%%%%%%%%%%%%%%%%%%%%%%%%%%%
\subsection{Ultraviolet and Infrared Cutoffs}
%%%%%%%%%%%%%%%%%%%%%%%%%%%%%%%%%%%%%%%%%%%%%%%%%%%%%%%%%%%%%%%%%%%%%%%%%%%%%%%%
Every finite representation introduces frequency limits
\[\omega_{\min},\qquad\omega_{\max}.\]
These cutoffs have direct physical consequences.
The infrared cutoff determines the longest temporal correlations present
in the stochastic field.
The ultraviolet cutoff limits the shortest fluctuations and suppresses
rapid stochastic kicks.
Because the zero-point spectrum satisfies
\[u(\omega)\propto\omega^3,\]
the total energy above frequency $\Omega$ grows as
\begin{equation}
\int^\Omega\omega^3\,d\omega\propto\Omega^4.
\end{equation}
Therefore,
high-frequency truncation removes a substantial fraction of the total
vacuum energy.
Whether this omission materially influences atomic dynamics remains one
of the principal open numerical questions.
%%%%%%%%%%%%%%%%%%%%%%%%%%%%%%%%%%%%%%%%%%%%%%%%%%%%%%%%%%%%%%%%%%%%%%%%%%%%%%%%
\subsection{Choice of Spectral Sampling}
%%%%%%%%%%%%%%%%%%%%%%%%%%%%%%%%%%%%%%%%%%%%%%%%%%%%%%%%%%%%%%%%%%%%%%%%%%%%%%%%
The frequencies
\[\{\omega_n\}\]
may be selected in several different ways.
Typical choices include
\begin{itemize}
\item
uniform spacing,
\item
logarithmic spacing,
\item
importance sampling,
\item
random frequency selection,
\item
quasi-random low-discrepancy sequences.
\end{itemize}
Each strategy introduces different statistical properties.
Uniform sampling provides systematic convergence but becomes expensive
at high frequencies.
Logarithmic sampling better resolves multiple frequency decades.
Importance sampling concentrates computational effort where the electron
responds most strongly.
Determining the optimal sampling strategy remains an active numerical
problem.
%%%%%%%%%%%%%%%%%%%%%%%%%%%%%%%%%%%%%%%%%%%%%%%%%%%%%%%%%%%%%%%%%%%%%%%%%%%%%%%%
\subsection{Random Phases}
%%%%%%%%%%%%%%%%%%%%%%%%%%%%%%%%%%%%%%%%%%%%%%%%%%%%%%%%%%%%%%%%%%%%%%%%%%%%%%%%
The stochastic character of the zero-point field is generated through
the random phases
\[\phi_n.\]
Ideally,
\[\phi_n\sim U(0,2\pi),\]
independently for every mode.
Poor-quality pseudorandom generators may introduce long-range
correlations that imitate physical effects.
Consequently,
modern simulations should employ statistically tested generators such
as
\begin{itemize}
\item
Mersenne Twister,
\item
PCG,
\item
xoshiro,
\item
counter-based generators suitable for parallel computation.
\end{itemize}
Independent simulations using different random generators provide an
important verification procedure.
%%%%%%%%%%%%%%%%%%%%%%%%%%%%%%%%%%%%%%%%%%%%%%%%%%%%%%%%%%%%%%%%%%%%%%%%%%%%%%%%
\subsection*{Critical Assessment}
%%%%%%%%%%%%%%%%%%%%%%%%%%%%%%%%%%%%%%%%%%%%%%%%%%%%%%%%%%%%%%%%%%%%%%%%%%%%%%%%
The numerical representation of the stochastic vacuum is not merely a
technical detail. It defines the stochastic process driving the electron
dynamics and therefore influences every computed observable.
For deterministic ordinary differential equations, refinement of the
time step alone is usually sufficient to establish convergence. In
stochastic electrodynamics this is not the case. Spectral refinement,
random-number quality, frequency sampling, and cutoff dependence all
alter the effective forcing experienced by the particle.
Accordingly, convincing numerical evidence requires simultaneous
convergence with respect to all of these parameters rather than any
single numerical approximation.
%%%%%%%%%%%%%%%%%%%%%%%%%%%%%%%%%%%%%%%%%%%%%%%%%%%%%%%%%%%%%%%%%%%%%%%%%%%%%%%%
\section{Time Integration Algorithms for SED}
\label{sec:TimeIntegration}
%%%%%%%%%%%%%%%%%%%%%%%%%%%%%%%%%%%%%%%%%%%%%%%%%%%%%%%%%%%%%%%%%%%%%%%%%%%%%%%%
The equations of stochastic electrodynamics are considerably more
difficult to integrate numerically than either ordinary Hamiltonian
systems or standard stochastic differential equations.
Three distinct features contribute to this complexity.
First,
the Coulomb force produces highly nonlinear motion, particularly during
close approaches to the nucleus.
Second,
the stochastic zero-point field contains a broad range of frequencies,
extending many orders of magnitude above the characteristic orbital
frequency.
Third,
radiation reaction introduces weak dissipation together with additional
high-order derivatives or their reduced-order approximations.
The resulting system is neither purely Hamiltonian nor a standard
Markovian stochastic differential equation.
Consequently, numerical methods developed for any one of these classes
cannot automatically be assumed appropriate.
%%%%%%%%%%%%%%%%%%%%%%%%%%%%%%%%%%%%%%%%%%%%%%%%%%%%%%%%%%%%%%%%%%%%%%%%%%%%%%%%
\subsection{Mathematical Structure}
%%%%%%%%%%%%%%%%%%%%%%%%%%%%%%%%%%%%%%%%%%%%%%%%%%%%%%%%%%%%%%%%%%%%%%%%%%%%%%%%
Ignoring radiation reaction for the moment, the equations may be written
\begin{align}
\dot{\mathbf r}&=\frac{\mathbf p}{m},\\
\dot{\mathbf p}&=-\nabla V(\mathbf r)+e\mathbf E_{\rm ZP}(t).
\end{align}
Unlike Langevin equations,
the stochastic force
\[\mathbf E_{\rm ZP}(t)\]
is a differentiable random function possessing finite temporal
correlations.
Thus the equations remain ordinary differential equations with random
coefficients rather than stochastic differential equations driven by
white noise.
This distinction is important because standard Itô and Stratonovich
integration methods are generally unnecessary unless additional white
noise approximations are introduced.
%%%%%%%%%%%%%%%%%%%%%%%%%%%%%%%%%%%%%%%%%%%%%%%%%%%%%%%%%%%%%%%%%%%%%%%%%%%%%%%%
\subsection{Characteristic Time Scales}
%%%%%%%%%%%%%%%%%%%%%%%%%%%%%%%%%%%%%%%%%%%%%%%%%%%%%%%%%%%%%%%%%%%%%%%%%%%%%%%%
A successful numerical algorithm must resolve several widely separated
time scales.
Among them are
\begin{itemize}
\item
the orbital period,
\item
the shortest represented vacuum wavelength,
\item
the radiation-reaction time
\[\tau_0=\frac{2e^2}{3mc^3},\]
\item
the longest vacuum correlation time,
\item
the total simulation time required to establish stationarity.
\end{itemize}
The ratio between the largest and smallest of these scales may exceed
many orders of magnitude.
The problem is therefore potentially stiff.
%%%%%%%%%%%%%%%%%%%%%%%%%%%%%%%%%%%%%%%%%%%%%%%%%%%%%%%%%%%%%%%%%%%%%%%%%%%%%%%%
\subsection{Explicit Runge--Kutta Methods}
%%%%%%%%%%%%%%%%%%%%%%%%%%%%%%%%%%%%%%%%%%%%%%%%%%%%%%%%%%%%%%%%%%%%%%%%%%%%%%%%
Most early SED simulations employed explicit Runge--Kutta schemes.
The advantages are well known.
\begin{itemize}
\item
Simple implementation.
\item
High local accuracy.
\item
Adaptive step-size control.
\item
Excellent performance for smooth trajectories.
\end{itemize}
However,
Runge--Kutta methods possess an important limitation.
They do not preserve the geometric structure of Hamiltonian dynamics.
Consequently,
small local errors may accumulate into significant long-term drift of
energy or angular momentum even in conservative systems.
%%%%%%%%%%%%%%%%%%%%%%%%%%%%%%%%%%%%%%%%%%%%%%%%%%%%%%%%%%%%%%%%%%%%%%%%%%%%%%%%
\subsection{Symplectic Integrators}
%%%%%%%%%%%%%%%%%%%%%%%%%%%%%%%%%%%%%%%%%%%%%%%%%%%%%%%%%%%%%%%%%%%%%%%%%%%%%%%%
For purely Hamiltonian systems,
symplectic methods preserve the canonical two-form,
\[\omega=dp_i\wedge dq^i.\]
Instead of exactly conserving energy,
they conserve a nearby modified Hamiltonian over exponentially long
times.
This remarkable property makes them the preferred methods for
\begin{itemize}
\item
planetary dynamics,
\item
molecular dynamics,
\item
accelerator physics,
\item
long-term celestial mechanics.
\end{itemize}
%%%%%%%%%%%%%%%%%%%%%%%%%%%%%%%%%%%%%%%%%%%%%%%%%%%%%%%%%%%%%%%%%%%%%%%%%%%%%%%%
\subsection{Can Symplectic Methods Be Used in SED?}
%%%%%%%%%%%%%%%%%%%%%%%%%%%%%%%%%%%%%%%%%%%%%%%%%%%%%%%%%%%%%%%%%%%%%%%%%%%%%%%%
The answer is subtle.
If radiation reaction is neglected,
the equations remain Hamiltonian with an explicitly time-dependent
forcing.
Symplectic integrators can then be generalized using extended phase
space or splitting methods.
Once radiation reaction is included,
strict symplecticity is lost because dissipation contracts phase-space
volume.
Nevertheless,
geometric integrators often continue to outperform conventional methods
by preserving the Hamiltonian part of the dynamics more faithfully.
%%%%%%%%%%%%%%%%%%%%%%%%%%%%%%%%%%%%%%%%%%%%%%%%%%%%%%%%%%%%%%%%%%%%%%%%%%%%%%%%
\subsection{Operator Splitting}
%%%%%%%%%%%%%%%%%%%%%%%%%%%%%%%%%%%%%%%%%%%%%%%%%%%%%%%%%%%%%%%%%%%%%%%%%%%%%%%%
A particularly attractive strategy is operator splitting.
The dynamics may be decomposed into three components,
\[L=L_{\rm Ham}+L_{\rm ZP}+L_{\rm RR},\]
corresponding respectively to
\begin{itemize}
\item
Hamiltonian motion,
\item
interaction with the stochastic field,
\item
radiation reaction.
\end{itemize}
Each component can be integrated separately,
after which the partial solutions are combined using symmetric
composition formulas,
for example
\[e^{hL}\approx e^{hL_{\rm RR}/2}e^{hL_{\rm ZP}/2}e^{hL_{\rm Ham}}
e^{hL_{\rm ZP}/2}e^{hL_{\rm RR}/2}.\]
Such methods preserve many geometric properties while incorporating weak
dissipation in a controlled manner.
%%%%%%%%%%%%%%%%%%%%%%%%%%%%%%%%%%%%%%%%%%%%%%%%%%%%%%%%%%%%%%%%%%%%%%%%%%%%%%%%
\subsection{Adaptive Time Stepping}
%%%%%%%%%%%%%%%%%%%%%%%%%%%%%%%%%%%%%%%%%%%%%%%%%%%%%%%%%%%%%%%%%%%%%%%%%%%%%%%%
Highly eccentric trajectories require locally refined time steps.
Near periapsis,
the Coulomb force varies rapidly,
whereas much larger steps become possible near aphelion.
Adaptive integration therefore offers substantial computational
advantages.
However,
changing the time step modifies the sampling of the stochastic field.
Interpolation of the zero-point field must therefore be performed with
care to avoid introducing artificial correlations.
%%%%%%%%%%%%%%%%%%%%%%%%%%%%%%%%%%%%%%%%%%%%%%%%%%%%%%%%%%%%%%%%%%%%%%%%%%%%%%%%
\subsection{Parallel Computation}
%%%%%%%%%%%%%%%%%%%%%%%%%%%%%%%%%%%%%%%%%%%%%%%%%%%%%%%%%%%%%%%%%%%%%%%%%%%%%%%%
The evaluation of thousands or millions of electromagnetic modes
dominates the computational cost.
Fortunately,
different modes contribute independently,
making the field evaluation naturally parallel.
Modern implementations may exploit
\begin{itemize}
\item
multicore processors,
\item
vector instructions,
\item
graphics processing units,
\item
distributed-memory clusters.
\end{itemize}
The stochastic trajectories themselves are likewise embarrassingly
parallel, since independent realizations require minimal communication.
%%%%%%%%%%%%%%%%%%%%%%%%%%%%%%%%%%%%%%%%%%%%%%%%%%%%%%%%%%%%%%%%%%%%%%%%%%%%%%%%
\subsection{Verification Strategy}
%%%%%%%%%%%%%%%%%%%%%%%%%%%%%%%%%%%%%%%%%%%%%%%%%%%%%%%%%%%%%%%%%%%%%%%%%%%%%%%%
Reliable numerical conclusions require more than agreement between two
runs.
At a minimum,
one should demonstrate convergence with respect to
\begin{enumerate}
\item
time-step refinement,
\item
spectral resolution,
\item
frequency cutoffs,
\item
integration algorithm,
\item
random-number generator,
\item
ensemble size,
\item
simulation duration.
\end{enumerate}
Only after all these tests have been passed can numerical agreement with
quantum mechanics be regarded as convincing.
%%%%%%%%%%%%%%%%%%%%%%%%%%%%%%%%%%%%%%%%%%%%%%%%%%%%%%%%%%%%%%%%%%%%%%%%%%%%%%%%
\subsection*{Critical Assessment}
%%%%%%%%%%%%%%%%%%%%%%%%%%%%%%%%%%%%%%%%%%%%%%%%%%%%%%%%%%%%%%%%%%%%%%%%%%%%%%%%
From the perspective of computational mechanics, stochastic
electrodynamics remains an exceptionally demanding numerical problem.
The coexistence of Hamiltonian motion, colored stochastic forcing, and
weak dissipation places it outside the scope of most standard numerical
analysis.
This observation also suggests a promising direction for future
research. During the past three decades, substantial advances have been
made in geometric numerical integration, structure-preserving
algorithms, exponential integrators, and high-performance stochastic
simulation. Relatively few of these developments have yet been applied
systematically to stochastic electrodynamics.
Consequently, it remains possible that improvements in numerical
methodology may prove as important as further analytical developments
for assessing the long-term stability of the SED hydrogen atom.
%%%%%%%%%%%%%%%%%%%%%%%%%%%%%%%%%%%%%%%%%%%%%%%%%%%%%%%%%%%%%%%%%%%%%%%%%%%%%%%%
\chapter{Stochastic Electrodynamics as an Open Hamiltonian System}
\label{ch:OpenHamiltonianSED}
%%%%%%%%%%%%%%%%%%%%%%%%%%%%%%%%%%%%%%%%%%%%%%%%%%%%%%%%%%%%%%%%%%%%%%%%%%%%%%%%
%%%%%%%%%%%%%%%%%%%%%%%%%%%%%%%%%%%%%%%%%%%%%%%%%%%%%%%%%%%%%%%%%%%%%%%%%%%%%%%%
\section{Introduction}
%%%%%%%%%%%%%%%%%%%%%%%%%%%%%%%%%%%%%%%%%%%%%%%%%%%%%%%%%%%%%%%%%%%%%%%%%%%%%%%%
The traditional presentation of stochastic electrodynamics introduces the
zero-point electromagnetic field as an externally prescribed stochastic
background acting upon charged particles. Although this viewpoint has
proved useful in practical calculations, it tends to obscure the fact
that the electromagnetic field is itself a dynamical system obeying
Maxwell's equations.
From a modern perspective it is more natural to regard stochastic
electrodynamics as an open Hamiltonian system obtained by eliminating
part of a larger deterministic dynamical system.
This point of view places stochastic electrodynamics within the general
framework of statistical mechanics, generalized Langevin equations,
projection-operator techniques, and open-system dynamics.
It also provides a direct conceptual bridge to stochastic Hamiltonian
mechanics and to the SHP framework developed later in this monograph.
%%%%%%%%%%%%%%%%%%%%%%%%%%%%%%%%%%%%%%%%%%%%%%%%%%%%%%%%%%%%%%%%%%%%%%%%%%%%%%%%
\section{The Complete Particle--Field System}
%%%%%%%%%%%%%%%%%%%%%%%%%%%%%%%%%%%%%%%%%%%%%%%%%%%%%%%%%%%%%%%%%%%%%%%%%%%%%%%%
The fundamental dynamical variables are
\[({\bf r},{\bf p}),\]
for the particle and
\[(A^\mu,\Pi^\mu)\]
for the electromagnetic field.
The complete state therefore belongs to the infinite-dimensional phase
space
\begin{equation}
\Gamma=\Gamma_{\rm particle}\times\Gamma_{\rm field}.
\end{equation}
The total Hamiltonian is
\begin{equation}
H=H_{\rm particle}+H_{\rm field}+H_{\rm interaction}.
\label{eq:TotalHamiltonian}
\end{equation}
Here
\begin{align}
H_{\rm particle}&=\frac{1}{2m}\left({\bf p}-\frac{e}{c}{\bf A}\right)^2+e\phi,\\
H_{\rm field}&=\frac{1}{8\pi}\int(E^2+B^2)\,d^3x,\\
H_{\rm interaction}&=-\frac{1}{c}\int J^\mu A_\mu \,d^3x.
\end{align}
No stochastic assumptions have yet been introduced.
The dynamics is completely deterministic.
%%%%%%%%%%%%%%%%%%%%%%%%%%%%%%%%%%%%%%%%%%%%%%%%%%%%%%%%%%%%%%%%%%%%%%%%%%%%%%%%
\section{Hamilton's Equations}
%%%%%%%%%%%%%%%%%%%%%%%%%%%%%%%%%%%%%%%%%%%%%%%%%%%%%%%%%%%%%%%%%%%%%%%%%%%%%%%%
The evolution is generated by Hamilton's equations,
\begin{align}
\dot q^i&=\frac{\partial H}{\partial p_i},\\
\dot p_i&=-\frac{\partial H}{\partial q^i},
\end{align}
together with the corresponding functional Hamilton equations for the
field variables,
\begin{align}
\dot A^\mu&=\frac{\delta H}{\delta\Pi_\mu},\\
\dot\Pi_\mu&=-\frac{\delta H}{\delta A^\mu}.
\end{align}
These equations are equivalent to Maxwell's equations coupled to the
Lorentz force.
The complete particle--field dynamics is therefore an ordinary
Hamiltonian flow on an infinite-dimensional phase space.
%%%%%%%%%%%%%%%%%%%%%%%%%%%%%%%%%%%%%%%%%%%%%%%%%%%%%%%%%%%%%%%%%%%%%%%%%%%%%%%%
\section{Projection onto Particle Variables}
%%%%%%%%%%%%%%%%%%%%%%%%%%%%%%%%%%%%%%%%%%%%%%%%%%%%%%%%%%%%%%%%%%%%%%%%%%%%%%%%
Suppose that only the particle variables
\[({\bf r},{\bf p})\]
are observed.
The field variables remain unobserved.
The observable state is therefore obtained by projecting
\[\Gamma\rightarrow\Gamma_{\rm particle}.\]
After projection,
the particle no longer evolves autonomously.
Instead,
its dynamics depends upon the unresolved field degrees of freedom.
The resulting equations are generally nonlocal in time and contain memory
effects.
%%%%%%%%%%%%%%%%%%%%%%%%%%%%%%%%%%%%%%%%%%%%%%%%%%%%%%%%%%%%%%%%%%%%%%%%%%%%%%%%
\section{Statistical Initial Conditions}
%%%%%%%%%%%%%%%%%%%%%%%%%%%%%%%%%%%%%%%%%%%%%%%%%%%%%%%%%%%%%%%%%%%%%%%%%%%%%%%%
Deterministic equations alone do not generate stochastic behavior.
Randomness enters through uncertainty concerning the initial state of
the electromagnetic field.
Instead of specifying
\[(A^\mu,\Pi^\mu)\]
exactly,
one assumes a probability measure
\[\mu(A,\Pi).\]
Particle observables are then obtained by averaging over this ensemble,
\begin{equation}
\langle O\rangle=\int O({\bf r},{\bf p};A,\Pi)\,d\mu(A,\Pi).
\end{equation}
Thus stochasticity originates in the ensemble of field initial
conditions rather than in the fundamental equations themselves.
%%%%%%%%%%%%%%%%%%%%%%%%%%%%%%%%%%%%%%%%%%%%%%%%%%%%%%%%%%%%%%%%%%%%%%%%%%%%%%%%
\section{Emergence of Memory}
%%%%%%%%%%%%%%%%%%%%%%%%%%%%%%%%%%%%%%%%%%%%%%%%%%%%%%%%%%%%%%%%%%%%%%%%%%%%%%%%
Eliminating the field variables generates an effective particle equation
\begin{equation}
m\ddot{\bf r}(t)={\bf F}_{\rm C}+\int_0^tK(t-s)\,{\bf r}(s)\,ds+{\boldsymbol\xi}(t),
\label{eq:GeneralizedLangevin}
\end{equation}
where
\[K(t)\]
is a memory kernel,
while
\[{\boldsymbol\xi}(t)\]
is the fluctuating force generated by the unknown initial field.
Equation
(\ref{eq:GeneralizedLangevin})
is a generalized Langevin equation.
Its structure is therefore identical to that encountered throughout
modern nonequilibrium statistical mechanics.
%%%%%%%%%%%%%%%%%%%%%%%%%%%%%%%%%%%%%%%%%%%%%%%%%%%%%%%%%%%%%%%%%%%%%%%%%%%%%%%%
\section{Fluctuation--Dissipation Relation}
%%%%%%%%%%%%%%%%%%%%%%%%%%%%%%%%%%%%%%%%%%%%%%%%%%%%%%%%%%%%%%%%%%%%%%%%%%%%%%%%
Because both the fluctuating force and the memory kernel originate from
the same eliminated field variables, they cannot be chosen
independently.
Their statistical properties satisfy a fluctuation--dissipation
relation,
\begin{equation}
\langle\xi_i(t)\xi_j(t')\rangle=k_BT K_{ij}(t-t')
\end{equation}
for thermal environments.
In stochastic electrodynamics the temperature is replaced by the
Lorentz-invariant zero-point spectrum.
Thus,
radiation reaction and vacuum fluctuations become two manifestations of
the same underlying electromagnetic field.
%%%%%%%%%%%%%%%%%%%%%%%%%%%%%%%%%%%%%%%%%%%%%%%%%%%%%%%%%%%%%%%%%%%%%%%%%%%%%%%%
\section{Conceptual Consequences}
%%%%%%%%%%%%%%%%%%%%%%%%%%%%%%%%%%%%%%%%%%%%%%%%%%%%%%%%%%%%%%%%%%%%%%%%%%%%%%%%
This reformulation changes the interpretation of stochastic
electrodynamics in several important respects.
\begin{enumerate}
\item
The complete theory remains deterministic.
\item
The stochastic equations describe only the reduced particle dynamics.
\item
Memory effects arise naturally.
\item
Radiation reaction and fluctuations possess a common origin.
\item
Invariant probability measures become the central mathematical objects.
\item
Modern methods of open-system theory become directly applicable.
\end{enumerate}
%%%%%%%%%%%%%%%%%%%%%%%%%%%%%%%%%%%%%%%%%%%%%%%%%%%%%%%%%%%%%%%%%%%%%%%%%%%%%%%%
\section*{Critical Assessment}
%%%%%%%%%%%%%%%%%%%%%%%%%%%%%%%%%%%%%%%%%%%%%%%%%%%%%%%%%%%%%%%%%%%%%%%%%%%%%%%%
Viewing stochastic electrodynamics as an open Hamiltonian system
clarifies many conceptual questions that have historically been treated
independently. In particular, it explains why stochastic forcing,
radiation reaction, and temporal memory appear simultaneously after the
electromagnetic field degrees of freedom are eliminated.
At the same time, this reformulation raises important mathematical
questions. The rigorous derivation of the effective particle dynamics
from the full Maxwell--Lorentz system remains highly nontrivial because
of self-interaction, gauge constraints, and the infinite number of field
degrees of freedom. Establishing the precise conditions under which the
reduced dynamics admits a generalized Langevin description is therefore
an important open problem.
Nevertheless, this perspective places stochastic electrodynamics within
the broader framework of modern nonequilibrium statistical mechanics,
where powerful analytical and numerical tools have been developed over
the past several decades. It also provides a natural bridge to the
stochastic Hamiltonian formulations discussed in the subsequent chapters
of this monograph.
